\section*{Lecture 11 - Platform-based Food Delivery}
\addcontentsline{toc}{section}{Lecture 11 - Platform-based Food Delivery}

\clearpage
\SlideHeader{Lecture 11 - Platform-based Food Delivery}{001}{Data Driven Decision Making}
\begin{minipage}[t]{0.405\textwidth}
\SlideImage{1}{../Materials/2026/3DM_11_Food-Delivery_SS26.pdf}
\begin{adjustbox}{max totalsize={\textwidth}{0.43\textheight},center}
\begin{minipage}{\textwidth}\scriptsize
\NoteSection{日本語訳}\begin{itemize}
\item データ駆動型意思決定
\item 講義 11 – Platform-based Food Delivery
\end{itemize}\NoteSection{試験対策ポイント}\begin{itemize}
\item Platform-based food deliveryは、需要・供給・推薦・配送の複数意思決定が連鎖する問題である。
\item food recommendationとdelivery optimizationを分けつつ、相互作用を説明する。
\end{itemize}
\end{minipage}
\end{adjustbox}
\end{minipage}\hfill
\begin{minipage}[t]{0.58\textwidth}
\begin{adjustbox}{max totalsize={\textwidth}{0.89\textheight},center}
\begin{minipage}{\textwidth}\scriptsize
\NoteSection{解説}このページでは「Data Driven Decision Making」を理解する。スライド上の表現は短いが、試験では定義、具体例、モデル上の意味、手法の限界まで説明できる必要がある。\par
Platform-based food delivery は、顧客、レストラン、配達員、プラットフォームが相互作用する多面的なサービスである。顧客は料理を選び、レストランは注文を調理し、配達員は複数注文を運び、プラットフォームは推薦、価格、割当、ルーティングを決める。\par
この問題では、単に最短ルートを作るだけでは不十分である。推薦システムが特定店舗に需要を集中させると、調理待ちや配達遅延が増える可能性がある。逆に、配送能力を考慮した推薦を行えば、需要を空いている店舗や配送しやすい地域へ誘導できる。\par
food recommendation では、content-based recommendation, collaborative filtering, ontology graph, diversificationなどが扱われる。顧客の好みだけでなく、料理の種類、価格、距離、調理時間、配達可能性などが推薦品質に影響する。\par
配送最適化では、注文受諾、配達員割当、ピックアップ順、ドロップオフ順、複数注文のバッチングが問題になる。時間窓、料理の品質劣化、顧客待ち時間、配達員の公平性なども考慮する必要がある。\par
試験では、食品配送をMDP/ADPとして設計する場合、状態に未処理注文、配達員位置、レストラン調理状態、需要予測を入れ、決定に推薦・注文受諾・割当・ルーティングを入れる。報酬には収益、遅延ペナルティ、顧客満足度、配達員稼働を含めるとよい。\par\NoteSection{確認問題と解答}\begin{enumerate}\item \textbf{L11-S001: 「Data Driven Decision Making」の概念を、定義・具体例・モデル上の意味の順に説明せよ。}\\定義としては、Platform-based food deliveryは、需要・供給・推薦・配送の複数意思決定が連鎖する問題である。。具体例では、配送・在庫・フードデリバリーのような現実問題を用い、状態、決定、外生情報、報酬、または情報モデル/意思決定モデルのどこに対応するかを明示する。モデル上の意味として、この概念が目的関数、制約、状態表現、将来価値、または方策選択にどのような影響を与えるかを書く。試験では、用語だけを列挙せず、入力データから意思決定、さらに現実の実装へつながる流れを文章で示すことが重要である。\item \textbf{L11-S001: このスライドの考え方を、講義全体のDDDMプロセスの中に位置づけよ。}\\DDDMプロセスでは、現実世界のデータを集約して情報モデルを作り、意思決定モデルまたは方策で行動を選び、実装後に結果を観測して次の状態へ進む。このスライドの概念は、そのどこか一部だけではなく、データ表現、意思決定、将来不確実性、計算可能性のいずれかに関わる。答案では、まずこの概念がどの段階に属するかを述べ、次に他の段階へどう影響するかを書く。例えば価値関数なら将来価値を現在決定へ戻す役割、FIFOなら旅行時間情報モデルの整合性を保証する役割、CFAなら目的関数を調整して将来に良い行動を誘導する役割である。\end{enumerate}
\end{minipage}
\end{adjustbox}
\end{minipage}


\clearpage
\SlideHeader{Lecture 11 - Platform-based Food Delivery}{002}{Content}
\begin{minipage}[t]{0.405\textwidth}
\SlideImage{2}{../Materials/2026/3DM_11_Food-Delivery_SS26.pdf}
\begin{adjustbox}{max totalsize={\textwidth}{0.43\textheight},center}
\begin{minipage}{\textwidth}\scriptsize
\NoteSection{日本語訳}\begin{itemize}
\item 内容
\item ▪ Platform-based Food Delivery（この用語は講義スライド上の専門語として扱う）
\item ▪ The Eater‘s Perspective（この用語は講義スライド上の専門語として扱う）
\item ▪ Uber Eat’s Perspective（この用語は講義スライド上の専門語として扱う）
\item ▪ The Delivery Perspective（この用語は講義スライド上の専門語として扱う）
\item ▪ The Restaurant‘s Perspective（この用語は講義スライド上の専門語として扱う）
\end{itemize}
\end{minipage}
\end{adjustbox}
\end{minipage}\hfill
\begin{minipage}[t]{0.58\textwidth}
\begin{adjustbox}{max totalsize={\textwidth}{0.89\textheight},center}
\begin{minipage}{\textwidth}\scriptsize
\NoteSection{注記}このページは表紙・目次・事務情報・導入画像が中心であるため、無理に確認問題や過去問を付けない。
\end{minipage}
\end{adjustbox}
\end{minipage}


\clearpage
\SlideHeader{Lecture 11 - Platform-based Food Delivery}{003}{Platform-based food delivery}
\begin{minipage}[t]{0.405\textwidth}
\SlideImage{3}{../Materials/2026/3DM_11_Food-Delivery_SS26.pdf}
\begin{adjustbox}{max totalsize={\textwidth}{0.43\textheight},center}
\begin{minipage}{\textwidth}\scriptsize
\NoteSection{日本語訳}\begin{itemize}
\item プラットフォーム型フードデリバリー
\item ▪ Platform-based business models act in a gig economy with volatile business partners（この用語は講義スライド上の専門語として扱う）
\item ▪ Uber Eats introduces the three sided-market place（この用語は講義スライド上の専門語として扱う）
\item ▪ Revenue mainly from restaurants (+user flat fee)（この用語は講義スライド上の専門語として扱う）
\item ▪ Payments are made to delivery drivers（この用語は講義スライド上の専門語として扱う）
\item ▪ Platform is not involved in physical transaction（この用語は講義スライド上の専門語として扱う）
\item https://www.intelivita.com/blog/food-delivery-business-model/（この用語は講義スライド上の専門語として扱う）
\end{itemize}\NoteSection{試験対策ポイント}\begin{itemize}
\item Platform-based food deliveryは、需要・供給・推薦・配送の複数意思決定が連鎖する問題である。
\item food recommendationとdelivery optimizationを分けつつ、相互作用を説明する。
\end{itemize}
\end{minipage}
\end{adjustbox}
\end{minipage}\hfill
\begin{minipage}[t]{0.58\textwidth}
\begin{adjustbox}{max totalsize={\textwidth}{0.89\textheight},center}
\begin{minipage}{\textwidth}\scriptsize
\NoteSection{解説}このページでは「Platform-based food delivery」を理解する。スライド上の表現は短いが、試験では定義、具体例、モデル上の意味、手法の限界まで説明できる必要がある。\par
Platform-based food delivery は、顧客、レストラン、配達員、プラットフォームが相互作用する多面的なサービスである。顧客は料理を選び、レストランは注文を調理し、配達員は複数注文を運び、プラットフォームは推薦、価格、割当、ルーティングを決める。\par
この問題では、単に最短ルートを作るだけでは不十分である。推薦システムが特定店舗に需要を集中させると、調理待ちや配達遅延が増える可能性がある。逆に、配送能力を考慮した推薦を行えば、需要を空いている店舗や配送しやすい地域へ誘導できる。\par
food recommendation では、content-based recommendation, collaborative filtering, ontology graph, diversificationなどが扱われる。顧客の好みだけでなく、料理の種類、価格、距離、調理時間、配達可能性などが推薦品質に影響する。\par
配送最適化では、注文受諾、配達員割当、ピックアップ順、ドロップオフ順、複数注文のバッチングが問題になる。時間窓、料理の品質劣化、顧客待ち時間、配達員の公平性なども考慮する必要がある。\par
試験では、食品配送をMDP/ADPとして設計する場合、状態に未処理注文、配達員位置、レストラン調理状態、需要予測を入れ、決定に推薦・注文受諾・割当・ルーティングを入れる。報酬には収益、遅延ペナルティ、顧客満足度、配達員稼働を含めるとよい。\par\NoteSection{確認問題と解答}\begin{enumerate}\item \textbf{L11-S003: 「Platform-based food delivery」の概念を、定義・具体例・モデル上の意味の順に説明せよ。}\\定義としては、Platform-based food deliveryは、需要・供給・推薦・配送の複数意思決定が連鎖する問題である。。具体例では、配送・在庫・フードデリバリーのような現実問題を用い、状態、決定、外生情報、報酬、または情報モデル/意思決定モデルのどこに対応するかを明示する。モデル上の意味として、この概念が目的関数、制約、状態表現、将来価値、または方策選択にどのような影響を与えるかを書く。試験では、用語だけを列挙せず、入力データから意思決定、さらに現実の実装へつながる流れを文章で示すことが重要である。\item \textbf{L11-S003: このスライドの考え方を、講義全体のDDDMプロセスの中に位置づけよ。}\\DDDMプロセスでは、現実世界のデータを集約して情報モデルを作り、意思決定モデルまたは方策で行動を選び、実装後に結果を観測して次の状態へ進む。このスライドの概念は、そのどこか一部だけではなく、データ表現、意思決定、将来不確実性、計算可能性のいずれかに関わる。答案では、まずこの概念がどの段階に属するかを述べ、次に他の段階へどう影響するかを書く。例えば価値関数なら将来価値を現在決定へ戻す役割、FIFOなら旅行時間情報モデルの整合性を保証する役割、CFAなら目的関数を調整して将来に良い行動を誘導する役割である。\end{enumerate}
\end{minipage}
\end{adjustbox}
\end{minipage}


\clearpage
\SlideHeader{Lecture 11 - Platform-based Food Delivery}{004}{Properties}
\begin{minipage}[t]{0.405\textwidth}
\SlideImage{4}{../Materials/2026/3DM_11_Food-Delivery_SS26.pdf}
\begin{adjustbox}{max totalsize={\textwidth}{0.43\textheight},center}
\begin{minipage}{\textwidth}\scriptsize
\NoteSection{日本語訳}\begin{itemize}
\item 性質
\item ▪ The platform business model is applicable in different regions and scales in any order（この用語は講義スライド上の専門語として扱う）
\item ▪ 決定s needed for the transactional control are taken in a fully automated fashion
\item ▪ Partners in the three-sided market place occur and vanish at random（この用語は講義スライド上の専門語として扱う）
\item ▪ The credibility of partners is maintained by monitoring, data analysis and ranking（この用語は講義スライド上の専門語として扱う）
\item ▪ Data driven 決定 making is the key feature for running such a platform
\end{itemize}\NoteSection{試験対策ポイント}\begin{itemize}
\item Platform-based food deliveryは、需要・供給・推薦・配送の複数意思決定が連鎖する問題である。
\item food recommendationとdelivery optimizationを分けつつ、相互作用を説明する。
\end{itemize}
\end{minipage}
\end{adjustbox}
\end{minipage}\hfill
\begin{minipage}[t]{0.58\textwidth}
\begin{adjustbox}{max totalsize={\textwidth}{0.89\textheight},center}
\begin{minipage}{\textwidth}\scriptsize
\NoteSection{解説}このページでは「Properties」を理解する。スライド上の表現は短いが、試験では定義、具体例、モデル上の意味、手法の限界まで説明できる必要がある。\par
Platform-based food delivery は、顧客、レストラン、配達員、プラットフォームが相互作用する多面的なサービスである。顧客は料理を選び、レストランは注文を調理し、配達員は複数注文を運び、プラットフォームは推薦、価格、割当、ルーティングを決める。\par
この問題では、単に最短ルートを作るだけでは不十分である。推薦システムが特定店舗に需要を集中させると、調理待ちや配達遅延が増える可能性がある。逆に、配送能力を考慮した推薦を行えば、需要を空いている店舗や配送しやすい地域へ誘導できる。\par
food recommendation では、content-based recommendation, collaborative filtering, ontology graph, diversificationなどが扱われる。顧客の好みだけでなく、料理の種類、価格、距離、調理時間、配達可能性などが推薦品質に影響する。\par
配送最適化では、注文受諾、配達員割当、ピックアップ順、ドロップオフ順、複数注文のバッチングが問題になる。時間窓、料理の品質劣化、顧客待ち時間、配達員の公平性なども考慮する必要がある。\par
試験では、食品配送をMDP/ADPとして設計する場合、状態に未処理注文、配達員位置、レストラン調理状態、需要予測を入れ、決定に推薦・注文受諾・割当・ルーティングを入れる。報酬には収益、遅延ペナルティ、顧客満足度、配達員稼働を含めるとよい。\par\NoteSection{確認問題と解答}\begin{enumerate}\item \textbf{L11-S004: 「Properties」の概念を、定義・具体例・モデル上の意味の順に説明せよ。}\\定義としては、Platform-based food deliveryは、需要・供給・推薦・配送の複数意思決定が連鎖する問題である。。具体例では、配送・在庫・フードデリバリーのような現実問題を用い、状態、決定、外生情報、報酬、または情報モデル/意思決定モデルのどこに対応するかを明示する。モデル上の意味として、この概念が目的関数、制約、状態表現、将来価値、または方策選択にどのような影響を与えるかを書く。試験では、用語だけを列挙せず、入力データから意思決定、さらに現実の実装へつながる流れを文章で示すことが重要である。\item \textbf{L11-S004: このスライドの考え方を、講義全体のDDDMプロセスの中に位置づけよ。}\\DDDMプロセスでは、現実世界のデータを集約して情報モデルを作り、意思決定モデルまたは方策で行動を選び、実装後に結果を観測して次の状態へ進む。このスライドの概念は、そのどこか一部だけではなく、データ表現、意思決定、将来不確実性、計算可能性のいずれかに関わる。答案では、まずこの概念がどの段階に属するかを述べ、次に他の段階へどう影響するかを書く。例えば価値関数なら将来価値を現在決定へ戻す役割、FIFOなら旅行時間情報モデルの整合性を保証する役割、CFAなら目的関数を調整して将来に良い行動を誘導する役割である。\end{enumerate}
\end{minipage}
\end{adjustbox}
\end{minipage}


\clearpage
\SlideHeader{Lecture 11 - Platform-based Food Delivery}{005}{Interaction of partner}
\begin{minipage}[t]{0.405\textwidth}
\SlideImage{5}{../Materials/2026/3DM_11_Food-Delivery_SS26.pdf}
\begin{adjustbox}{max totalsize={\textwidth}{0.43\textheight},center}
\begin{minipage}{\textwidth}\scriptsize
\NoteSection{日本語訳}\begin{itemize}
\item 参加者間の相互作用
\item https://www.infoq.com/fr/articles/uber-eats-time-predictions/（この用語は講義スライド上の専門語として扱う）
\end{itemize}\NoteSection{試験対策ポイント}\begin{itemize}
\item Platform-based food deliveryは、需要・供給・推薦・配送の複数意思決定が連鎖する問題である。
\item food recommendationとdelivery optimizationを分けつつ、相互作用を説明する。
\end{itemize}
\end{minipage}
\end{adjustbox}
\end{minipage}\hfill
\begin{minipage}[t]{0.58\textwidth}
\begin{adjustbox}{max totalsize={\textwidth}{0.89\textheight},center}
\begin{minipage}{\textwidth}\scriptsize
\NoteSection{解説}このページでは「Interaction of partner」を理解する。スライド上の表現は短いが、試験では定義、具体例、モデル上の意味、手法の限界まで説明できる必要がある。\par
Platform-based food delivery は、顧客、レストラン、配達員、プラットフォームが相互作用する多面的なサービスである。顧客は料理を選び、レストランは注文を調理し、配達員は複数注文を運び、プラットフォームは推薦、価格、割当、ルーティングを決める。\par
この問題では、単に最短ルートを作るだけでは不十分である。推薦システムが特定店舗に需要を集中させると、調理待ちや配達遅延が増える可能性がある。逆に、配送能力を考慮した推薦を行えば、需要を空いている店舗や配送しやすい地域へ誘導できる。\par
food recommendation では、content-based recommendation, collaborative filtering, ontology graph, diversificationなどが扱われる。顧客の好みだけでなく、料理の種類、価格、距離、調理時間、配達可能性などが推薦品質に影響する。\par
配送最適化では、注文受諾、配達員割当、ピックアップ順、ドロップオフ順、複数注文のバッチングが問題になる。時間窓、料理の品質劣化、顧客待ち時間、配達員の公平性なども考慮する必要がある。\par
試験では、食品配送をMDP/ADPとして設計する場合、状態に未処理注文、配達員位置、レストラン調理状態、需要予測を入れ、決定に推薦・注文受諾・割当・ルーティングを入れる。報酬には収益、遅延ペナルティ、顧客満足度、配達員稼働を含めるとよい。\par\NoteSection{確認問題と解答}\begin{enumerate}\item \textbf{L11-S005: 「Interaction of partner」の概念を、定義・具体例・モデル上の意味の順に説明せよ。}\\定義としては、Platform-based food deliveryは、需要・供給・推薦・配送の複数意思決定が連鎖する問題である。。具体例では、配送・在庫・フードデリバリーのような現実問題を用い、状態、決定、外生情報、報酬、または情報モデル/意思決定モデルのどこに対応するかを明示する。モデル上の意味として、この概念が目的関数、制約、状態表現、将来価値、または方策選択にどのような影響を与えるかを書く。試験では、用語だけを列挙せず、入力データから意思決定、さらに現実の実装へつながる流れを文章で示すことが重要である。\item \textbf{L11-S005: このスライドの考え方を、講義全体のDDDMプロセスの中に位置づけよ。}\\DDDMプロセスでは、現実世界のデータを集約して情報モデルを作り、意思決定モデルまたは方策で行動を選び、実装後に結果を観測して次の状態へ進む。このスライドの概念は、そのどこか一部だけではなく、データ表現、意思決定、将来不確実性、計算可能性のいずれかに関わる。答案では、まずこの概念がどの段階に属するかを述べ、次に他の段階へどう影響するかを書く。例えば価値関数なら将来価値を現在決定へ戻す役割、FIFOなら旅行時間情報モデルの整合性を保証する役割、CFAなら目的関数を調整して将来に良い行動を誘導する役割である。\end{enumerate}
\end{minipage}
\end{adjustbox}
\end{minipage}


\clearpage
\SlideHeader{Lecture 11 - Platform-based Food Delivery}{006}{Content}
\begin{minipage}[t]{0.405\textwidth}
\SlideImage{6}{../Materials/2026/3DM_11_Food-Delivery_SS26.pdf}
\begin{adjustbox}{max totalsize={\textwidth}{0.43\textheight},center}
\begin{minipage}{\textwidth}\scriptsize
\NoteSection{日本語訳}\begin{itemize}
\item 内容
\item ▪ Platform-based Food Delivery（この用語は講義スライド上の専門語として扱う）
\item ▪ The Eater‘s Perspective（この用語は講義スライド上の専門語として扱う）
\item ▪ Uber Eat’s Perspective（この用語は講義スライド上の専門語として扱う）
\item ▪ The Delivery Perspective（この用語は講義スライド上の専門語として扱う）
\item ▪ The Restaurant‘s Perspective（この用語は講義スライド上の専門語として扱う）
\end{itemize}
\end{minipage}
\end{adjustbox}
\end{minipage}\hfill
\begin{minipage}[t]{0.58\textwidth}
\begin{adjustbox}{max totalsize={\textwidth}{0.89\textheight},center}
\begin{minipage}{\textwidth}\scriptsize
\NoteSection{注記}このページは表紙・目次・事務情報・導入画像が中心であるため、無理に確認問題や過去問を付けない。
\end{minipage}
\end{adjustbox}
\end{minipage}


\clearpage
\SlideHeader{Lecture 11 - Platform-based Food Delivery}{007}{Stages in generation of food recommendation}
\begin{minipage}[t]{0.405\textwidth}
\SlideImage{7}{../Materials/2026/3DM_11_Food-Delivery_SS26.pdf}
\begin{adjustbox}{max totalsize={\textwidth}{0.43\textheight},center}
\begin{minipage}{\textwidth}\scriptsize
\NoteSection{日本語訳}\begin{itemize}
\item Stages in generation of food 推薦
\item 1. Non-personalized 推薦 focus on Content
\item very few restaurants only based（この用語は講義スライド上の専門語として扱う）
\item 2. Content-based 推薦s
\item have led to a meal over-specialization（この用語は講義スライド上の専門語として扱う）
\item Collaborative（この用語は講義スライド上の専門語として扱う）
\item 3. Collaborative filtering is used for filtering（この用語は講義スライド上の専門語として扱う）
\item personalized 推薦s
\item ➢ 3. requires the clustering of eaters in（この用語は講義スライド上の専門語として扱う）
\end{itemize}\NoteSection{試験対策ポイント}\begin{itemize}
\item Platform-based food deliveryは、需要・供給・推薦・配送の複数意思決定が連鎖する問題である。
\item food recommendationとdelivery optimizationを分けつつ、相互作用を説明する。
\end{itemize}
\end{minipage}
\end{adjustbox}
\end{minipage}\hfill
\begin{minipage}[t]{0.58\textwidth}
\begin{adjustbox}{max totalsize={\textwidth}{0.89\textheight},center}
\begin{minipage}{\textwidth}\scriptsize
\NoteSection{解説}このページでは「Stages in generation of food recommendation」を理解する。スライド上の表現は短いが、試験では定義、具体例、モデル上の意味、手法の限界まで説明できる必要がある。\par
Platform-based food delivery は、顧客、レストラン、配達員、プラットフォームが相互作用する多面的なサービスである。顧客は料理を選び、レストランは注文を調理し、配達員は複数注文を運び、プラットフォームは推薦、価格、割当、ルーティングを決める。\par
この問題では、単に最短ルートを作るだけでは不十分である。推薦システムが特定店舗に需要を集中させると、調理待ちや配達遅延が増える可能性がある。逆に、配送能力を考慮した推薦を行えば、需要を空いている店舗や配送しやすい地域へ誘導できる。\par
food recommendation では、content-based recommendation, collaborative filtering, ontology graph, diversificationなどが扱われる。顧客の好みだけでなく、料理の種類、価格、距離、調理時間、配達可能性などが推薦品質に影響する。\par
配送最適化では、注文受諾、配達員割当、ピックアップ順、ドロップオフ順、複数注文のバッチングが問題になる。時間窓、料理の品質劣化、顧客待ち時間、配達員の公平性なども考慮する必要がある。\par
試験では、食品配送をMDP/ADPとして設計する場合、状態に未処理注文、配達員位置、レストラン調理状態、需要予測を入れ、決定に推薦・注文受諾・割当・ルーティングを入れる。報酬には収益、遅延ペナルティ、顧客満足度、配達員稼働を含めるとよい。\par\NoteSection{確認問題と解答}\begin{enumerate}\item \textbf{L11-S007: 「Stages in generation of food recommendation」の概念を、定義・具体例・モデル上の意味の順に説明せよ。}\\定義としては、Platform-based food deliveryは、需要・供給・推薦・配送の複数意思決定が連鎖する問題である。。具体例では、配送・在庫・フードデリバリーのような現実問題を用い、状態、決定、外生情報、報酬、または情報モデル/意思決定モデルのどこに対応するかを明示する。モデル上の意味として、この概念が目的関数、制約、状態表現、将来価値、または方策選択にどのような影響を与えるかを書く。試験では、用語だけを列挙せず、入力データから意思決定、さらに現実の実装へつながる流れを文章で示すことが重要である。\item \textbf{L11-S007: このスライドの考え方を、講義全体のDDDMプロセスの中に位置づけよ。}\\DDDMプロセスでは、現実世界のデータを集約して情報モデルを作り、意思決定モデルまたは方策で行動を選び、実装後に結果を観測して次の状態へ進む。このスライドの概念は、そのどこか一部だけではなく、データ表現、意思決定、将来不確実性、計算可能性のいずれかに関わる。答案では、まずこの概念がどの段階に属するかを述べ、次に他の段階へどう影響するかを書く。例えば価値関数なら将来価値を現在決定へ戻す役割、FIFOなら旅行時間情報モデルの整合性を保証する役割、CFAなら目的関数を調整して将来に良い行動を誘導する役割である。\end{enumerate}
\end{minipage}
\end{adjustbox}
\end{minipage}


\clearpage
\SlideHeader{Lecture 11 - Platform-based Food Delivery}{008}{The eater‘s perspective (non-personalized recommendation)}
\begin{minipage}[t]{0.405\textwidth}
\SlideImage{8}{../Materials/2026/3DM_11_Food-Delivery_SS26.pdf}
\begin{adjustbox}{max totalsize={\textwidth}{0.43\textheight},center}
\begin{minipage}{\textwidth}\scriptsize
\NoteSection{日本語訳}\begin{itemize}
\item The eater‘s perspective (non-personalized 推薦)
\item ▪ The eater has got a preferred cuisine（この用語は講義スライド上の専門語として扱う）
\item ▪ She prefers food from high ranked restaurants（この用語は講義スライド上の専門語として扱う）
\item ▪ Fast delivery is of utmost importance（この用語は講義スライド上の専門語として扱う）
\item ▪ Reasonable pricing is desirable（この用語は講義スライド上の専門語として扱う）
\item ▪ However, restaurant choices may be too similar（この用語は講義スライド上の専門語として扱う）
\item ▪ There should be some range of food offered to the eater（この用語は講義スライド上の専門語として扱う）
\item ▪ Diverse food offerings maintain eater‘s retention（この用語は講義スライド上の専門語として扱う）
\item https://www.uber.com/en-DE/blog/uber-eats-trip-（この用語は講義スライド上の専門語として扱う）
\end{itemize}\NoteSection{試験対策ポイント}\begin{itemize}
\item Platform-based food deliveryは、需要・供給・推薦・配送の複数意思決定が連鎖する問題である。
\item food recommendationとdelivery optimizationを分けつつ、相互作用を説明する。
\end{itemize}
\end{minipage}
\end{adjustbox}
\end{minipage}\hfill
\begin{minipage}[t]{0.58\textwidth}
\begin{adjustbox}{max totalsize={\textwidth}{0.89\textheight},center}
\begin{minipage}{\textwidth}\scriptsize
\NoteSection{解説}このページでは「The eater‘s perspective (non-personalized recommendation)」を理解する。スライド上の表現は短いが、試験では定義、具体例、モデル上の意味、手法の限界まで説明できる必要がある。\par
Platform-based food delivery は、顧客、レストラン、配達員、プラットフォームが相互作用する多面的なサービスである。顧客は料理を選び、レストランは注文を調理し、配達員は複数注文を運び、プラットフォームは推薦、価格、割当、ルーティングを決める。\par
この問題では、単に最短ルートを作るだけでは不十分である。推薦システムが特定店舗に需要を集中させると、調理待ちや配達遅延が増える可能性がある。逆に、配送能力を考慮した推薦を行えば、需要を空いている店舗や配送しやすい地域へ誘導できる。\par
food recommendation では、content-based recommendation, collaborative filtering, ontology graph, diversificationなどが扱われる。顧客の好みだけでなく、料理の種類、価格、距離、調理時間、配達可能性などが推薦品質に影響する。\par
配送最適化では、注文受諾、配達員割当、ピックアップ順、ドロップオフ順、複数注文のバッチングが問題になる。時間窓、料理の品質劣化、顧客待ち時間、配達員の公平性なども考慮する必要がある。\par
試験では、食品配送をMDP/ADPとして設計する場合、状態に未処理注文、配達員位置、レストラン調理状態、需要予測を入れ、決定に推薦・注文受諾・割当・ルーティングを入れる。報酬には収益、遅延ペナルティ、顧客満足度、配達員稼働を含めるとよい。\par\NoteSection{確認問題と解答}\begin{enumerate}\item \textbf{L11-S008: 「The eater‘s perspective (non-personalized recommendation)」の概念を、定義・具体例・モデル上の意味の順に説明せよ。}\\定義としては、Platform-based food deliveryは、需要・供給・推薦・配送の複数意思決定が連鎖する問題である。。具体例では、配送・在庫・フードデリバリーのような現実問題を用い、状態、決定、外生情報、報酬、または情報モデル/意思決定モデルのどこに対応するかを明示する。モデル上の意味として、この概念が目的関数、制約、状態表現、将来価値、または方策選択にどのような影響を与えるかを書く。試験では、用語だけを列挙せず、入力データから意思決定、さらに現実の実装へつながる流れを文章で示すことが重要である。\item \textbf{L11-S008: このスライドの考え方を、講義全体のDDDMプロセスの中に位置づけよ。}\\DDDMプロセスでは、現実世界のデータを集約して情報モデルを作り、意思決定モデルまたは方策で行動を選び、実装後に結果を観測して次の状態へ進む。このスライドの概念は、そのどこか一部だけではなく、データ表現、意思決定、将来不確実性、計算可能性のいずれかに関わる。答案では、まずこの概念がどの段階に属するかを述べ、次に他の段階へどう影響するかを書く。例えば価値関数なら将来価値を現在決定へ戻す役割、FIFOなら旅行時間情報モデルの整合性を保証する役割、CFAなら目的関数を調整して将来に良い行動を誘導する役割である。\end{enumerate}
\end{minipage}
\end{adjustbox}
\end{minipage}


\clearpage
\SlideHeader{Lecture 11 - Platform-based Food Delivery}{009}{Food ontology graph (content based recommendation)}
\begin{minipage}[t]{0.405\textwidth}
\SlideImage{9}{../Materials/2026/3DM_11_Food-Delivery_SS26.pdf}
\begin{adjustbox}{max totalsize={\textwidth}{0.43\textheight},center}
\begin{minipage}{\textwidth}\scriptsize
\NoteSection{日本語訳}\begin{itemize}
\item 食のオントロジーグラフ（内容ベース推薦）
\item https://www.uber.com/en-DE/blog/uber-eats-query-understanding/（この用語は講義スライド上の専門語として扱う）
\end{itemize}\NoteSection{試験対策ポイント}\begin{itemize}
\item Platform-based food deliveryは、需要・供給・推薦・配送の複数意思決定が連鎖する問題である。
\item food recommendationとdelivery optimizationを分けつつ、相互作用を説明する。
\end{itemize}
\end{minipage}
\end{adjustbox}
\end{minipage}\hfill
\begin{minipage}[t]{0.58\textwidth}
\begin{adjustbox}{max totalsize={\textwidth}{0.89\textheight},center}
\begin{minipage}{\textwidth}\scriptsize
\NoteSection{解説}このページでは「Food ontology graph (content based recommendation)」を理解する。スライド上の表現は短いが、試験では定義、具体例、モデル上の意味、手法の限界まで説明できる必要がある。\par
Platform-based food delivery は、顧客、レストラン、配達員、プラットフォームが相互作用する多面的なサービスである。顧客は料理を選び、レストランは注文を調理し、配達員は複数注文を運び、プラットフォームは推薦、価格、割当、ルーティングを決める。\par
この問題では、単に最短ルートを作るだけでは不十分である。推薦システムが特定店舗に需要を集中させると、調理待ちや配達遅延が増える可能性がある。逆に、配送能力を考慮した推薦を行えば、需要を空いている店舗や配送しやすい地域へ誘導できる。\par
food recommendation では、content-based recommendation, collaborative filtering, ontology graph, diversificationなどが扱われる。顧客の好みだけでなく、料理の種類、価格、距離、調理時間、配達可能性などが推薦品質に影響する。\par
配送最適化では、注文受諾、配達員割当、ピックアップ順、ドロップオフ順、複数注文のバッチングが問題になる。時間窓、料理の品質劣化、顧客待ち時間、配達員の公平性なども考慮する必要がある。\par
試験では、食品配送をMDP/ADPとして設計する場合、状態に未処理注文、配達員位置、レストラン調理状態、需要予測を入れ、決定に推薦・注文受諾・割当・ルーティングを入れる。報酬には収益、遅延ペナルティ、顧客満足度、配達員稼働を含めるとよい。\par\NoteSection{確認問題と解答}\begin{enumerate}\item \textbf{L11-S009: 「Food ontology graph (content based recommendation)」の概念を、定義・具体例・モデル上の意味の順に説明せよ。}\\定義としては、Platform-based food deliveryは、需要・供給・推薦・配送の複数意思決定が連鎖する問題である。。具体例では、配送・在庫・フードデリバリーのような現実問題を用い、状態、決定、外生情報、報酬、または情報モデル/意思決定モデルのどこに対応するかを明示する。モデル上の意味として、この概念が目的関数、制約、状態表現、将来価値、または方策選択にどのような影響を与えるかを書く。試験では、用語だけを列挙せず、入力データから意思決定、さらに現実の実装へつながる流れを文章で示すことが重要である。\item \textbf{L11-S009: このスライドの考え方を、講義全体のDDDMプロセスの中に位置づけよ。}\\DDDMプロセスでは、現実世界のデータを集約して情報モデルを作り、意思決定モデルまたは方策で行動を選び、実装後に結果を観測して次の状態へ進む。このスライドの概念は、そのどこか一部だけではなく、データ表現、意思決定、将来不確実性、計算可能性のいずれかに関わる。答案では、まずこの概念がどの段階に属するかを述べ、次に他の段階へどう影響するかを書く。例えば価値関数なら将来価値を現在決定へ戻す役割、FIFOなら旅行時間情報モデルの整合性を保証する役割、CFAなら目的関数を調整して将来に良い行動を誘導する役割である。\end{enumerate}
\end{minipage}
\end{adjustbox}
\end{minipage}


\clearpage
\SlideHeader{Lecture 11 - Platform-based Food Delivery}{010}{Content based diversification of meal offerings}
\begin{minipage}[t]{0.405\textwidth}
\SlideImage{10}{../Materials/2026/3DM_11_Food-Delivery_SS26.pdf}
\begin{adjustbox}{max totalsize={\textwidth}{0.43\textheight},center}
\begin{minipage}{\textwidth}\scriptsize
\NoteSection{日本語訳}\begin{itemize}
\item Content based diversification of meal offerings（この用語は講義スライド上の専門語として扱う）
\item ▪ Ontology subgraph links eater query with alternative cuisine（この用語は講義スライド上の専門語として扱う）
\item ▪ Having searched for Soba, also Udon may be offered（この用語は講義スライド上の専門語として扱う）
\item ▪ Broaden the scope may also offer Sushi, or Chinese dishes（この用語は講義スライド上の専門語として扱う）
\item ▪ Ontologies are built offline by incorporating expert knowledge（この用語は講義スライド上の専門語として扱う）
\item ▪ The online counterpart for diversification is learning from orders（この用語は講義スライド上の専門語として扱う）
\item https://www.uber.com/en-DE/blog/uber-eats-query-understanding/（この用語は講義スライド上の専門語として扱う）
\end{itemize}\NoteSection{試験対策ポイント}\begin{itemize}
\item Platform-based food deliveryは、需要・供給・推薦・配送の複数意思決定が連鎖する問題である。
\item food recommendationとdelivery optimizationを分けつつ、相互作用を説明する。
\end{itemize}
\end{minipage}
\end{adjustbox}
\end{minipage}\hfill
\begin{minipage}[t]{0.58\textwidth}
\begin{adjustbox}{max totalsize={\textwidth}{0.89\textheight},center}
\begin{minipage}{\textwidth}\scriptsize
\NoteSection{解説}このページでは「Content based diversification of meal offerings」を理解する。スライド上の表現は短いが、試験では定義、具体例、モデル上の意味、手法の限界まで説明できる必要がある。\par
Platform-based food delivery は、顧客、レストラン、配達員、プラットフォームが相互作用する多面的なサービスである。顧客は料理を選び、レストランは注文を調理し、配達員は複数注文を運び、プラットフォームは推薦、価格、割当、ルーティングを決める。\par
この問題では、単に最短ルートを作るだけでは不十分である。推薦システムが特定店舗に需要を集中させると、調理待ちや配達遅延が増える可能性がある。逆に、配送能力を考慮した推薦を行えば、需要を空いている店舗や配送しやすい地域へ誘導できる。\par
food recommendation では、content-based recommendation, collaborative filtering, ontology graph, diversificationなどが扱われる。顧客の好みだけでなく、料理の種類、価格、距離、調理時間、配達可能性などが推薦品質に影響する。\par
配送最適化では、注文受諾、配達員割当、ピックアップ順、ドロップオフ順、複数注文のバッチングが問題になる。時間窓、料理の品質劣化、顧客待ち時間、配達員の公平性なども考慮する必要がある。\par
試験では、食品配送をMDP/ADPとして設計する場合、状態に未処理注文、配達員位置、レストラン調理状態、需要予測を入れ、決定に推薦・注文受諾・割当・ルーティングを入れる。報酬には収益、遅延ペナルティ、顧客満足度、配達員稼働を含めるとよい。\par\NoteSection{確認問題と解答}\begin{enumerate}\item \textbf{L11-S010: 「Content based diversification of meal offerings」の概念を、定義・具体例・モデル上の意味の順に説明せよ。}\\定義としては、Platform-based food deliveryは、需要・供給・推薦・配送の複数意思決定が連鎖する問題である。。具体例では、配送・在庫・フードデリバリーのような現実問題を用い、状態、決定、外生情報、報酬、または情報モデル/意思決定モデルのどこに対応するかを明示する。モデル上の意味として、この概念が目的関数、制約、状態表現、将来価値、または方策選択にどのような影響を与えるかを書く。試験では、用語だけを列挙せず、入力データから意思決定、さらに現実の実装へつながる流れを文章で示すことが重要である。\item \textbf{L11-S010: このスライドの考え方を、講義全体のDDDMプロセスの中に位置づけよ。}\\DDDMプロセスでは、現実世界のデータを集約して情報モデルを作り、意思決定モデルまたは方策で行動を選び、実装後に結果を観測して次の状態へ進む。このスライドの概念は、そのどこか一部だけではなく、データ表現、意思決定、将来不確実性、計算可能性のいずれかに関わる。答案では、まずこの概念がどの段階に属するかを述べ、次に他の段階へどう影響するかを書く。例えば価値関数なら将来価値を現在決定へ戻す役割、FIFOなら旅行時間情報モデルの整合性を保証する役割、CFAなら目的関数を調整して将来に良い行動を誘導する役割である。\end{enumerate}
\end{minipage}
\end{adjustbox}
\end{minipage}


\clearpage
\SlideHeader{Lecture 11 - Platform-based Food Delivery}{011}{Online diversification (collaborative filtering)}
\begin{minipage}[t]{0.405\textwidth}
\SlideImage{11}{../Materials/2026/3DM_11_Food-Delivery_SS26.pdf}
\begin{adjustbox}{max totalsize={\textwidth}{0.43\textheight},center}
\begin{minipage}{\textwidth}\scriptsize
\NoteSection{日本語訳}\begin{itemize}
\item Online diversification (協調フィルタリング)
\item ▪ Queries are tied to restaurants by historic orders（この用語は講義スライド上の専門語として扱う）
\item ▪ Queries referring to the same restaurant define a context（この用語は講義スライド上の専門語として扱う）
\item ▪ The context is used later on for query expansion（この用語は講義スライド上の専門語として扱う）
\item https://www.uber.com/en-DE/blog/uber-eats-query-understanding/（この用語は講義スライド上の専門語として扱う）
\item ▪ For Tan Tan Noodle request, restaurants are recommended that appear in related contexts（この用語は講義スライド上の専門語として扱う）
\end{itemize}\NoteSection{試験対策ポイント}\begin{itemize}
\item Platform-based food deliveryは、需要・供給・推薦・配送の複数意思決定が連鎖する問題である。
\item food recommendationとdelivery optimizationを分けつつ、相互作用を説明する。
\end{itemize}
\end{minipage}
\end{adjustbox}
\end{minipage}\hfill
\begin{minipage}[t]{0.58\textwidth}
\begin{adjustbox}{max totalsize={\textwidth}{0.89\textheight},center}
\begin{minipage}{\textwidth}\scriptsize
\NoteSection{解説}このページでは「Online diversification (collaborative filtering)」を理解する。スライド上の表現は短いが、試験では定義、具体例、モデル上の意味、手法の限界まで説明できる必要がある。\par
Platform-based food delivery は、顧客、レストラン、配達員、プラットフォームが相互作用する多面的なサービスである。顧客は料理を選び、レストランは注文を調理し、配達員は複数注文を運び、プラットフォームは推薦、価格、割当、ルーティングを決める。\par
この問題では、単に最短ルートを作るだけでは不十分である。推薦システムが特定店舗に需要を集中させると、調理待ちや配達遅延が増える可能性がある。逆に、配送能力を考慮した推薦を行えば、需要を空いている店舗や配送しやすい地域へ誘導できる。\par
food recommendation では、content-based recommendation, collaborative filtering, ontology graph, diversificationなどが扱われる。顧客の好みだけでなく、料理の種類、価格、距離、調理時間、配達可能性などが推薦品質に影響する。\par
配送最適化では、注文受諾、配達員割当、ピックアップ順、ドロップオフ順、複数注文のバッチングが問題になる。時間窓、料理の品質劣化、顧客待ち時間、配達員の公平性なども考慮する必要がある。\par
試験では、食品配送をMDP/ADPとして設計する場合、状態に未処理注文、配達員位置、レストラン調理状態、需要予測を入れ、決定に推薦・注文受諾・割当・ルーティングを入れる。報酬には収益、遅延ペナルティ、顧客満足度、配達員稼働を含めるとよい。\par\NoteSection{確認問題と解答}\begin{enumerate}\item \textbf{L11-S011: 「Online diversification (collaborative filtering)」の概念を、定義・具体例・モデル上の意味の順に説明せよ。}\\定義としては、Platform-based food deliveryは、需要・供給・推薦・配送の複数意思決定が連鎖する問題である。。具体例では、配送・在庫・フードデリバリーのような現実問題を用い、状態、決定、外生情報、報酬、または情報モデル/意思決定モデルのどこに対応するかを明示する。モデル上の意味として、この概念が目的関数、制約、状態表現、将来価値、または方策選択にどのような影響を与えるかを書く。試験では、用語だけを列挙せず、入力データから意思決定、さらに現実の実装へつながる流れを文章で示すことが重要である。\item \textbf{L11-S011: このスライドの考え方を、講義全体のDDDMプロセスの中に位置づけよ。}\\DDDMプロセスでは、現実世界のデータを集約して情報モデルを作り、意思決定モデルまたは方策で行動を選び、実装後に結果を観測して次の状態へ進む。このスライドの概念は、そのどこか一部だけではなく、データ表現、意思決定、将来不確実性、計算可能性のいずれかに関わる。答案では、まずこの概念がどの段階に属するかを述べ、次に他の段階へどう影響するかを書く。例えば価値関数なら将来価値を現在決定へ戻す役割、FIFOなら旅行時間情報モデルの整合性を保証する役割、CFAなら目的関数を調整して将来に良い行動を誘導する役割である。\end{enumerate}
\end{minipage}
\end{adjustbox}
\end{minipage}


\clearpage
\SlideHeader{Lecture 11 - Platform-based Food Delivery}{012}{Content}
\begin{minipage}[t]{0.405\textwidth}
\SlideImage{12}{../Materials/2026/3DM_11_Food-Delivery_SS26.pdf}
\begin{adjustbox}{max totalsize={\textwidth}{0.43\textheight},center}
\begin{minipage}{\textwidth}\scriptsize
\NoteSection{日本語訳}\begin{itemize}
\item 内容
\item ▪ Platform-based Food Delivery（この用語は講義スライド上の専門語として扱う）
\item ▪ The Eater‘s Perspective（この用語は講義スライド上の専門語として扱う）
\item ▪ Uber Eat’s Perspective（この用語は講義スライド上の専門語として扱う）
\item ▪ The Delivery Perspective（この用語は講義スライド上の専門語として扱う）
\item ▪ The Restaurant‘s Perspective（この用語は講義スライド上の専門語として扱う）
\end{itemize}
\end{minipage}
\end{adjustbox}
\end{minipage}\hfill
\begin{minipage}[t]{0.58\textwidth}
\begin{adjustbox}{max totalsize={\textwidth}{0.89\textheight},center}
\begin{minipage}{\textwidth}\scriptsize
\NoteSection{注記}このページは表紙・目次・事務情報・導入画像が中心であるため、無理に確認問題や過去問を付けない。
\end{minipage}
\end{adjustbox}
\end{minipage}


\clearpage
\SlideHeader{Lecture 11 - Platform-based Food Delivery}{013}{Controlling the marketplace}
\begin{minipage}[t]{0.405\textwidth}
\SlideImage{13}{../Materials/2026/3DM_11_Food-Delivery_SS26.pdf}
\begin{adjustbox}{max totalsize={\textwidth}{0.43\textheight},center}
\begin{minipage}{\textwidth}\scriptsize
\NoteSection{日本語訳}\begin{itemize}
\item Controlling the marketplace（この用語は講義スライド上の専門語として扱う）
\item https://www.uber.com/en-DE/blog/uber-eats-trip-（この用語は講義スライド上の専門語として扱う）
\item optimization/（この用語は講義スライド上の専門語として扱う）
\end{itemize}\NoteSection{試験対策ポイント}\begin{itemize}
\item Platform-based food deliveryは、需要・供給・推薦・配送の複数意思決定が連鎖する問題である。
\item food recommendationとdelivery optimizationを分けつつ、相互作用を説明する。
\end{itemize}
\end{minipage}
\end{adjustbox}
\end{minipage}\hfill
\begin{minipage}[t]{0.58\textwidth}
\begin{adjustbox}{max totalsize={\textwidth}{0.89\textheight},center}
\begin{minipage}{\textwidth}\scriptsize
\NoteSection{解説}このページでは「Controlling the marketplace」を理解する。スライド上の表現は短いが、試験では定義、具体例、モデル上の意味、手法の限界まで説明できる必要がある。\par
Platform-based food delivery は、顧客、レストラン、配達員、プラットフォームが相互作用する多面的なサービスである。顧客は料理を選び、レストランは注文を調理し、配達員は複数注文を運び、プラットフォームは推薦、価格、割当、ルーティングを決める。\par
この問題では、単に最短ルートを作るだけでは不十分である。推薦システムが特定店舗に需要を集中させると、調理待ちや配達遅延が増える可能性がある。逆に、配送能力を考慮した推薦を行えば、需要を空いている店舗や配送しやすい地域へ誘導できる。\par
food recommendation では、content-based recommendation, collaborative filtering, ontology graph, diversificationなどが扱われる。顧客の好みだけでなく、料理の種類、価格、距離、調理時間、配達可能性などが推薦品質に影響する。\par
配送最適化では、注文受諾、配達員割当、ピックアップ順、ドロップオフ順、複数注文のバッチングが問題になる。時間窓、料理の品質劣化、顧客待ち時間、配達員の公平性なども考慮する必要がある。\par
試験では、食品配送をMDP/ADPとして設計する場合、状態に未処理注文、配達員位置、レストラン調理状態、需要予測を入れ、決定に推薦・注文受諾・割当・ルーティングを入れる。報酬には収益、遅延ペナルティ、顧客満足度、配達員稼働を含めるとよい。\par\NoteSection{確認問題と解答}\begin{enumerate}\item \textbf{L11-S013: 「Controlling the marketplace」の概念を、定義・具体例・モデル上の意味の順に説明せよ。}\\定義としては、Platform-based food deliveryは、需要・供給・推薦・配送の複数意思決定が連鎖する問題である。。具体例では、配送・在庫・フードデリバリーのような現実問題を用い、状態、決定、外生情報、報酬、または情報モデル/意思決定モデルのどこに対応するかを明示する。モデル上の意味として、この概念が目的関数、制約、状態表現、将来価値、または方策選択にどのような影響を与えるかを書く。試験では、用語だけを列挙せず、入力データから意思決定、さらに現実の実装へつながる流れを文章で示すことが重要である。\item \textbf{L11-S013: このスライドの考え方を、講義全体のDDDMプロセスの中に位置づけよ。}\\DDDMプロセスでは、現実世界のデータを集約して情報モデルを作り、意思決定モデルまたは方策で行動を選び、実装後に結果を観測して次の状態へ進む。このスライドの概念は、そのどこか一部だけではなく、データ表現、意思決定、将来不確実性、計算可能性のいずれかに関わる。答案では、まずこの概念がどの段階に属するかを述べ、次に他の段階へどう影響するかを書く。例えば価値関数なら将来価値を現在決定へ戻す役割、FIFOなら旅行時間情報モデルの整合性を保証する役割、CFAなら目的関数を調整して将来に良い行動を誘導する役割である。\end{enumerate}
\end{minipage}
\end{adjustbox}
\end{minipage}


\clearpage
\SlideHeader{Lecture 11 - Platform-based Food Delivery}{014}{Ranking of recommendations}
\begin{minipage}[t]{0.405\textwidth}
\SlideImage{14}{../Materials/2026/3DM_11_Food-Delivery_SS26.pdf}
\begin{adjustbox}{max totalsize={\textwidth}{0.43\textheight},center}
\begin{minipage}{\textwidth}\scriptsize
\NoteSection{日本語訳}\begin{itemize}
\item Ranking of 推薦s
\item ▪ Approximately 80\% of order are picked from recommended meals shown in app（この用語は講義スライド上の専門語として扱う）
\item ▪ Recommendations should be generated in a way that aligns with partner‘s objectives（この用語は講義スライド上の専門語として扱う）
\item Partner 目的
\item Conversion（この用語は講義スライド上の専門語として扱う）
\item Eater（この用語は講義スライド上の専門語として扱う）
\item Retention（この用語は講義スライド上の専門語として扱う）
\item Fairness（この用語は講義スライド上の専門語として扱う）
\item Restaurant（この用語は講義スライド上の専門語として扱う）
\end{itemize}\NoteSection{試験対策ポイント}\begin{itemize}
\item Platform-based food deliveryは、需要・供給・推薦・配送の複数意思決定が連鎖する問題である。
\item food recommendationとdelivery optimizationを分けつつ、相互作用を説明する。
\end{itemize}
\end{minipage}
\end{adjustbox}
\end{minipage}\hfill
\begin{minipage}[t]{0.58\textwidth}
\begin{adjustbox}{max totalsize={\textwidth}{0.89\textheight},center}
\begin{minipage}{\textwidth}\scriptsize
\NoteSection{解説}このページでは「Ranking of recommendations」を理解する。スライド上の表現は短いが、試験では定義、具体例、モデル上の意味、手法の限界まで説明できる必要がある。\par
Platform-based food delivery は、顧客、レストラン、配達員、プラットフォームが相互作用する多面的なサービスである。顧客は料理を選び、レストランは注文を調理し、配達員は複数注文を運び、プラットフォームは推薦、価格、割当、ルーティングを決める。\par
この問題では、単に最短ルートを作るだけでは不十分である。推薦システムが特定店舗に需要を集中させると、調理待ちや配達遅延が増える可能性がある。逆に、配送能力を考慮した推薦を行えば、需要を空いている店舗や配送しやすい地域へ誘導できる。\par
food recommendation では、content-based recommendation, collaborative filtering, ontology graph, diversificationなどが扱われる。顧客の好みだけでなく、料理の種類、価格、距離、調理時間、配達可能性などが推薦品質に影響する。\par
配送最適化では、注文受諾、配達員割当、ピックアップ順、ドロップオフ順、複数注文のバッチングが問題になる。時間窓、料理の品質劣化、顧客待ち時間、配達員の公平性なども考慮する必要がある。\par
試験では、食品配送をMDP/ADPとして設計する場合、状態に未処理注文、配達員位置、レストラン調理状態、需要予測を入れ、決定に推薦・注文受諾・割当・ルーティングを入れる。報酬には収益、遅延ペナルティ、顧客満足度、配達員稼働を含めるとよい。\par\NoteSection{確認問題と解答}\begin{enumerate}\item \textbf{L11-S014: 「Ranking of recommendations」の概念を、定義・具体例・モデル上の意味の順に説明せよ。}\\定義としては、Platform-based food deliveryは、需要・供給・推薦・配送の複数意思決定が連鎖する問題である。。具体例では、配送・在庫・フードデリバリーのような現実問題を用い、状態、決定、外生情報、報酬、または情報モデル/意思決定モデルのどこに対応するかを明示する。モデル上の意味として、この概念が目的関数、制約、状態表現、将来価値、または方策選択にどのような影響を与えるかを書く。試験では、用語だけを列挙せず、入力データから意思決定、さらに現実の実装へつながる流れを文章で示すことが重要である。\item \textbf{L11-S014: このスライドの考え方を、講義全体のDDDMプロセスの中に位置づけよ。}\\DDDMプロセスでは、現実世界のデータを集約して情報モデルを作り、意思決定モデルまたは方策で行動を選び、実装後に結果を観測して次の状態へ進む。このスライドの概念は、そのどこか一部だけではなく、データ表現、意思決定、将来不確実性、計算可能性のいずれかに関わる。答案では、まずこの概念がどの段階に属するかを述べ、次に他の段階へどう影響するかを書く。例えば価値関数なら将来価値を現在決定へ戻す役割、FIFOなら旅行時間情報モデルの整合性を保証する役割、CFAなら目的関数を調整して将来に良い行動を誘導する役割である。\end{enumerate}
\end{minipage}
\end{adjustbox}
\end{minipage}


\clearpage
\SlideHeader{Lecture 11 - Platform-based Food Delivery}{015}{Incorporate multiple (contradicting) objectives}
\begin{minipage}[t]{0.405\textwidth}
\SlideImage{15}{../Materials/2026/3DM_11_Food-Delivery_SS26.pdf}
\begin{adjustbox}{max totalsize={\textwidth}{0.43\textheight},center}
\begin{minipage}{\textwidth}\scriptsize
\NoteSection{日本語訳}\begin{itemize}
\item Incorporate multiple (contradicting) objectives（この用語は講義スライド上の専門語として扱う）
\item ▪ Gross booking value is defined as the full amount of payments processed on the platform. It is（この用語は講義スライド上の専門語として扱う）
\item achieved by placing pricy orders at popular restaurants and addresses restaurant earnings（この用語は講義スライド上の専門語として扱う）
\item ▪ Marketplace fairness is in the interest of new or less popular restaurants. Earnings from such（この用語は講義スライド上の専門語として扱う）
\item restaurants should be generated in order to prevent such restaurants from leaving the platform（この用語は講義スライド上の専門語として扱う）
\item https://www.adjust.com/glossary/conversion-rate/（この用語は講義スライド上の専門語として扱う）
\item ▪ A conversion rate records the percentage of users who have completed a desired action. Here,（この用語は講義スライド上の専門語として扱う）
\item its defines a user interaction that results in a meal order. This strengthen eater retention（この用語は講義スライド上の専門語として扱う）
\item ▪ Gross booking value conflicts with marketplace fairness（この用語は講義スライド上の専門語として扱う）
\end{itemize}\NoteSection{試験対策ポイント}\begin{itemize}
\item Platform-based food deliveryは、需要・供給・推薦・配送の複数意思決定が連鎖する問題である。
\item food recommendationとdelivery optimizationを分けつつ、相互作用を説明する。
\end{itemize}
\end{minipage}
\end{adjustbox}
\end{minipage}\hfill
\begin{minipage}[t]{0.58\textwidth}
\begin{adjustbox}{max totalsize={\textwidth}{0.89\textheight},center}
\begin{minipage}{\textwidth}\scriptsize
\NoteSection{解説}このページでは「Incorporate multiple (contradicting) objectives」を理解する。スライド上の表現は短いが、試験では定義、具体例、モデル上の意味、手法の限界まで説明できる必要がある。\par
Platform-based food delivery は、顧客、レストラン、配達員、プラットフォームが相互作用する多面的なサービスである。顧客は料理を選び、レストランは注文を調理し、配達員は複数注文を運び、プラットフォームは推薦、価格、割当、ルーティングを決める。\par
この問題では、単に最短ルートを作るだけでは不十分である。推薦システムが特定店舗に需要を集中させると、調理待ちや配達遅延が増える可能性がある。逆に、配送能力を考慮した推薦を行えば、需要を空いている店舗や配送しやすい地域へ誘導できる。\par
food recommendation では、content-based recommendation, collaborative filtering, ontology graph, diversificationなどが扱われる。顧客の好みだけでなく、料理の種類、価格、距離、調理時間、配達可能性などが推薦品質に影響する。\par
配送最適化では、注文受諾、配達員割当、ピックアップ順、ドロップオフ順、複数注文のバッチングが問題になる。時間窓、料理の品質劣化、顧客待ち時間、配達員の公平性なども考慮する必要がある。\par
試験では、食品配送をMDP/ADPとして設計する場合、状態に未処理注文、配達員位置、レストラン調理状態、需要予測を入れ、決定に推薦・注文受諾・割当・ルーティングを入れる。報酬には収益、遅延ペナルティ、顧客満足度、配達員稼働を含めるとよい。\par\NoteSection{確認問題と解答}\begin{enumerate}\item \textbf{L11-S015: 「Incorporate multiple (contradicting) objectives」の概念を、定義・具体例・モデル上の意味の順に説明せよ。}\\定義としては、Platform-based food deliveryは、需要・供給・推薦・配送の複数意思決定が連鎖する問題である。。具体例では、配送・在庫・フードデリバリーのような現実問題を用い、状態、決定、外生情報、報酬、または情報モデル/意思決定モデルのどこに対応するかを明示する。モデル上の意味として、この概念が目的関数、制約、状態表現、将来価値、または方策選択にどのような影響を与えるかを書く。試験では、用語だけを列挙せず、入力データから意思決定、さらに現実の実装へつながる流れを文章で示すことが重要である。\item \textbf{L11-S015: このスライドの考え方を、講義全体のDDDMプロセスの中に位置づけよ。}\\DDDMプロセスでは、現実世界のデータを集約して情報モデルを作り、意思決定モデルまたは方策で行動を選び、実装後に結果を観測して次の状態へ進む。このスライドの概念は、そのどこか一部だけではなく、データ表現、意思決定、将来不確実性、計算可能性のいずれかに関わる。答案では、まずこの概念がどの段階に属するかを述べ、次に他の段階へどう影響するかを書く。例えば価値関数なら将来価値を現在決定へ戻す役割、FIFOなら旅行時間情報モデルの整合性を保証する役割、CFAなら目的関数を調整して将来に良い行動を誘導する役割である。\end{enumerate}
\end{minipage}
\end{adjustbox}
\end{minipage}


\clearpage
\SlideHeader{Lecture 11 - Platform-based Food Delivery}{016}{Modeling (1)}
\begin{minipage}[t]{0.405\textwidth}
\SlideImage{16}{../Materials/2026/3DM_11_Food-Delivery_SS26.pdf}
\begin{adjustbox}{max totalsize={\textwidth}{0.43\textheight},center}
\begin{minipage}{\textwidth}\scriptsize
\NoteSection{日本語訳}\begin{itemize}
\item Modeling (1)（この用語は講義スライド上の専門語として扱う）
\item Model decides about the priority of recommending (reasonable) restaurants to eater（この用語は講義スライド上の専門語として扱う）
\item 𝑥𝑖𝑗 = 決定 variable that represents the probability of recommending restaurant 𝑗 to eater 𝑖
\item 𝐼𝑖𝑗 = indicator of restaurant 𝑗 being recommended by eater 𝑖（この用語は講義スライド上の専門語として扱う）
\item 𝑂𝑖𝑗 = indicator of eater 𝑖 ordering from restaurant 𝑗（この用語は講義スライド上の専門語として扱う）
\item 𝐺𝑖𝑗 = dollar amount of the order between eater 𝑖 and restaurant 𝑗（この用語は講義スライド上の専門語として扱う）
\item Conversion rate 𝑇 𝑋 = 𝐸[σ𝑖 σ𝑗 𝐼𝑖𝑗 𝑂𝑖𝑗 ] = σ𝑖 σ𝑗 𝑥𝑖𝑗 𝑝𝑖𝑗 with 𝑝𝑖𝑗 predicted conversion rate（この用語は講義スライド上の専門語として扱う）
\item Gross booking 𝐺 𝑋 = 𝐸[σ𝑖 σ𝑗 𝐼𝑖𝑗 𝑂𝑖𝑗 𝐺𝑖𝑗 ] = σ𝑖 σ𝑗 𝑥𝑖𝑗 𝑝𝑖𝑗 𝑔𝑖𝑗 with 𝑔𝑖𝑗 期待される booking value
\item Both, 𝑝𝑖𝑗 = 𝐸[𝑂𝑖𝑗 ] and 𝑔𝑖𝑗 = 𝐸[𝐺𝑖𝑗 ] are personalized machine learning models that use（この用語は講義スライド上の専門語として扱う）
\end{itemize}\NoteSection{試験対策ポイント}\begin{itemize}
\item Platform-based food deliveryは、需要・供給・推薦・配送の複数意思決定が連鎖する問題である。
\item food recommendationとdelivery optimizationを分けつつ、相互作用を説明する。
\end{itemize}
\end{minipage}
\end{adjustbox}
\end{minipage}\hfill
\begin{minipage}[t]{0.58\textwidth}
\begin{adjustbox}{max totalsize={\textwidth}{0.89\textheight},center}
\begin{minipage}{\textwidth}\scriptsize
\NoteSection{解説}このページでは「Modeling (1)」を理解する。スライド上の表現は短いが、試験では定義、具体例、モデル上の意味、手法の限界まで説明できる必要がある。\par
Platform-based food delivery は、顧客、レストラン、配達員、プラットフォームが相互作用する多面的なサービスである。顧客は料理を選び、レストランは注文を調理し、配達員は複数注文を運び、プラットフォームは推薦、価格、割当、ルーティングを決める。\par
この問題では、単に最短ルートを作るだけでは不十分である。推薦システムが特定店舗に需要を集中させると、調理待ちや配達遅延が増える可能性がある。逆に、配送能力を考慮した推薦を行えば、需要を空いている店舗や配送しやすい地域へ誘導できる。\par
food recommendation では、content-based recommendation, collaborative filtering, ontology graph, diversificationなどが扱われる。顧客の好みだけでなく、料理の種類、価格、距離、調理時間、配達可能性などが推薦品質に影響する。\par
配送最適化では、注文受諾、配達員割当、ピックアップ順、ドロップオフ順、複数注文のバッチングが問題になる。時間窓、料理の品質劣化、顧客待ち時間、配達員の公平性なども考慮する必要がある。\par
試験では、食品配送をMDP/ADPとして設計する場合、状態に未処理注文、配達員位置、レストラン調理状態、需要予測を入れ、決定に推薦・注文受諾・割当・ルーティングを入れる。報酬には収益、遅延ペナルティ、顧客満足度、配達員稼働を含めるとよい。\par\NoteSection{確認問題と解答}\begin{enumerate}\item \textbf{L11-S016: 「Modeling (1)」の概念を、定義・具体例・モデル上の意味の順に説明せよ。}\\定義としては、Platform-based food deliveryは、需要・供給・推薦・配送の複数意思決定が連鎖する問題である。。具体例では、配送・在庫・フードデリバリーのような現実問題を用い、状態、決定、外生情報、報酬、または情報モデル/意思決定モデルのどこに対応するかを明示する。モデル上の意味として、この概念が目的関数、制約、状態表現、将来価値、または方策選択にどのような影響を与えるかを書く。試験では、用語だけを列挙せず、入力データから意思決定、さらに現実の実装へつながる流れを文章で示すことが重要である。\item \textbf{L11-S016: このスライドの考え方を、講義全体のDDDMプロセスの中に位置づけよ。}\\DDDMプロセスでは、現実世界のデータを集約して情報モデルを作り、意思決定モデルまたは方策で行動を選び、実装後に結果を観測して次の状態へ進む。このスライドの概念は、そのどこか一部だけではなく、データ表現、意思決定、将来不確実性、計算可能性のいずれかに関わる。答案では、まずこの概念がどの段階に属するかを述べ、次に他の段階へどう影響するかを書く。例えば価値関数なら将来価値を現在決定へ戻す役割、FIFOなら旅行時間情報モデルの整合性を保証する役割、CFAなら目的関数を調整して将来に良い行動を誘導する役割である。\end{enumerate}
\end{minipage}
\end{adjustbox}
\end{minipage}


\clearpage
\SlideHeader{Lecture 11 - Platform-based Food Delivery}{017}{Modeling (2)}
\begin{minipage}[t]{0.405\textwidth}
\SlideImage{17}{../Materials/2026/3DM_11_Food-Delivery_SS26.pdf}
\begin{adjustbox}{max totalsize={\textwidth}{0.43\textheight},center}
\begin{minipage}{\textwidth}\scriptsize
\NoteSection{日本語訳}\begin{itemize}
\item Modeling (2)（この用語は講義スライド上の専門語として扱う）
\item 𝑈 = (𝑢𝑖𝑗 ) equal probability for all restaurants 𝑗 to be recommended to eater 𝑖 with σ𝑗 𝑢𝑖𝑗 = 1, ∀𝑖（この用語は講義スライド上の専門語として扱う）
\item 𝑋 ⋆ = probabilities of optimal solution that can achieve best possible conversion rate 𝑇(𝑋 ⋆ )（この用語は講義スライド上の専門語として扱う）
\item Parameter (0 < 𝛼 < 1) for required conversion rate and 𝜆 for quadratic fairness penalty（この用語は講義スライド上の専門語として扱う）
\item Maximize gross booking value max 𝐺 𝑋 − 𝜆 𝑋 − 𝑈 2（この用語は講義スライド上の専門語として扱う）
\item 𝑋 2
\item 𝑠. 𝑡. 𝑇 𝑋 ≥ 𝛼 𝑇(𝑋 ⋆ )
\item 𝑥
\item 𝑢 |𝑥 − 𝑢| ෍ 𝑥𝑖𝑗 = 1, ∀𝑖
\end{itemize}\NoteSection{試験対策ポイント}\begin{itemize}
\item Platform-based food deliveryは、需要・供給・推薦・配送の複数意思決定が連鎖する問題である。
\item food recommendationとdelivery optimizationを分けつつ、相互作用を説明する。
\end{itemize}
\end{minipage}
\end{adjustbox}
\end{minipage}\hfill
\begin{minipage}[t]{0.58\textwidth}
\begin{adjustbox}{max totalsize={\textwidth}{0.89\textheight},center}
\begin{minipage}{\textwidth}\scriptsize
\NoteSection{解説}このページでは「Modeling (2)」を理解する。スライド上の表現は短いが、試験では定義、具体例、モデル上の意味、手法の限界まで説明できる必要がある。\par
Platform-based food delivery は、顧客、レストラン、配達員、プラットフォームが相互作用する多面的なサービスである。顧客は料理を選び、レストランは注文を調理し、配達員は複数注文を運び、プラットフォームは推薦、価格、割当、ルーティングを決める。\par
この問題では、単に最短ルートを作るだけでは不十分である。推薦システムが特定店舗に需要を集中させると、調理待ちや配達遅延が増える可能性がある。逆に、配送能力を考慮した推薦を行えば、需要を空いている店舗や配送しやすい地域へ誘導できる。\par
food recommendation では、content-based recommendation, collaborative filtering, ontology graph, diversificationなどが扱われる。顧客の好みだけでなく、料理の種類、価格、距離、調理時間、配達可能性などが推薦品質に影響する。\par
配送最適化では、注文受諾、配達員割当、ピックアップ順、ドロップオフ順、複数注文のバッチングが問題になる。時間窓、料理の品質劣化、顧客待ち時間、配達員の公平性なども考慮する必要がある。\par
試験では、食品配送をMDP/ADPとして設計する場合、状態に未処理注文、配達員位置、レストラン調理状態、需要予測を入れ、決定に推薦・注文受諾・割当・ルーティングを入れる。報酬には収益、遅延ペナルティ、顧客満足度、配達員稼働を含めるとよい。\par\NoteSection{確認問題と解答}\begin{enumerate}\item \textbf{L11-S017: 「Modeling (2)」の概念を、定義・具体例・モデル上の意味の順に説明せよ。}\\定義としては、Platform-based food deliveryは、需要・供給・推薦・配送の複数意思決定が連鎖する問題である。。具体例では、配送・在庫・フードデリバリーのような現実問題を用い、状態、決定、外生情報、報酬、または情報モデル/意思決定モデルのどこに対応するかを明示する。モデル上の意味として、この概念が目的関数、制約、状態表現、将来価値、または方策選択にどのような影響を与えるかを書く。試験では、用語だけを列挙せず、入力データから意思決定、さらに現実の実装へつながる流れを文章で示すことが重要である。\item \textbf{L11-S017: このスライドの考え方を、講義全体のDDDMプロセスの中に位置づけよ。}\\DDDMプロセスでは、現実世界のデータを集約して情報モデルを作り、意思決定モデルまたは方策で行動を選び、実装後に結果を観測して次の状態へ進む。このスライドの概念は、そのどこか一部だけではなく、データ表現、意思決定、将来不確実性、計算可能性のいずれかに関わる。答案では、まずこの概念がどの段階に属するかを述べ、次に他の段階へどう影響するかを書く。例えば価値関数なら将来価値を現在決定へ戻す役割、FIFOなら旅行時間情報モデルの整合性を保証する役割、CFAなら目的関数を調整して将来に良い行動を誘導する役割である。\end{enumerate}
\end{minipage}
\end{adjustbox}
\end{minipage}


\clearpage
\SlideHeader{Lecture 11 - Platform-based Food Delivery}{018}{Modeling (3)}
\begin{minipage}[t]{0.405\textwidth}
\SlideImage{18}{../Materials/2026/3DM_11_Food-Delivery_SS26.pdf}
\begin{adjustbox}{max totalsize={\textwidth}{0.43\textheight},center}
\begin{minipage}{\textwidth}\scriptsize
\NoteSection{日本語訳}\begin{itemize}
\item Modeling (3)（この用語は講義スライド上の専門語として扱う）
\item ▪ Parameter λ penalizes gross booking value that results from restaurant 推薦s
\item ▪ A small λ favors recommended restaurants and drives out low performing restaurants（この用語は講義スライド上の専門語として扱う）
\item ▪ A high λ enforces that gross booking value is generated from low performers only（この用語は講義スライド上の専門語として扱う）
\item ▪ Parameter α enforces a high conversion rate.（この用語は講義スライド上の専門語として扱う）
\item ▪ Setting α ≈ 1 discharges promising solutions（この用語は講義スライド上の専門語として扱う）
\item ▪ Setting α ≈ 0 allows for a poor conversion rate（この用語は講義スライド上の専門語として扱う）
\item ▪ Both parameters have to be tuned online（この用語は講義スライド上の専門語として扱う）
\item ▪ A multi-armed bandit maintains different parameter sets in parallel（この用語は講義スライド上の専門語として扱う）
\end{itemize}\NoteSection{試験対策ポイント}\begin{itemize}
\item Platform-based food deliveryは、需要・供給・推薦・配送の複数意思決定が連鎖する問題である。
\item food recommendationとdelivery optimizationを分けつつ、相互作用を説明する。
\end{itemize}
\end{minipage}
\end{adjustbox}
\end{minipage}\hfill
\begin{minipage}[t]{0.58\textwidth}
\begin{adjustbox}{max totalsize={\textwidth}{0.89\textheight},center}
\begin{minipage}{\textwidth}\scriptsize
\NoteSection{解説}このページでは「Modeling (3)」を理解する。スライド上の表現は短いが、試験では定義、具体例、モデル上の意味、手法の限界まで説明できる必要がある。\par
Platform-based food delivery は、顧客、レストラン、配達員、プラットフォームが相互作用する多面的なサービスである。顧客は料理を選び、レストランは注文を調理し、配達員は複数注文を運び、プラットフォームは推薦、価格、割当、ルーティングを決める。\par
この問題では、単に最短ルートを作るだけでは不十分である。推薦システムが特定店舗に需要を集中させると、調理待ちや配達遅延が増える可能性がある。逆に、配送能力を考慮した推薦を行えば、需要を空いている店舗や配送しやすい地域へ誘導できる。\par
food recommendation では、content-based recommendation, collaborative filtering, ontology graph, diversificationなどが扱われる。顧客の好みだけでなく、料理の種類、価格、距離、調理時間、配達可能性などが推薦品質に影響する。\par
配送最適化では、注文受諾、配達員割当、ピックアップ順、ドロップオフ順、複数注文のバッチングが問題になる。時間窓、料理の品質劣化、顧客待ち時間、配達員の公平性なども考慮する必要がある。\par
試験では、食品配送をMDP/ADPとして設計する場合、状態に未処理注文、配達員位置、レストラン調理状態、需要予測を入れ、決定に推薦・注文受諾・割当・ルーティングを入れる。報酬には収益、遅延ペナルティ、顧客満足度、配達員稼働を含めるとよい。\par\NoteSection{確認問題と解答}\begin{enumerate}\item \textbf{L11-S018: 「Modeling (3)」の概念を、定義・具体例・モデル上の意味の順に説明せよ。}\\定義としては、Platform-based food deliveryは、需要・供給・推薦・配送の複数意思決定が連鎖する問題である。。具体例では、配送・在庫・フードデリバリーのような現実問題を用い、状態、決定、外生情報、報酬、または情報モデル/意思決定モデルのどこに対応するかを明示する。モデル上の意味として、この概念が目的関数、制約、状態表現、将来価値、または方策選択にどのような影響を与えるかを書く。試験では、用語だけを列挙せず、入力データから意思決定、さらに現実の実装へつながる流れを文章で示すことが重要である。\item \textbf{L11-S018: このスライドの考え方を、講義全体のDDDMプロセスの中に位置づけよ。}\\DDDMプロセスでは、現実世界のデータを集約して情報モデルを作り、意思決定モデルまたは方策で行動を選び、実装後に結果を観測して次の状態へ進む。このスライドの概念は、そのどこか一部だけではなく、データ表現、意思決定、将来不確実性、計算可能性のいずれかに関わる。答案では、まずこの概念がどの段階に属するかを述べ、次に他の段階へどう影響するかを書く。例えば価値関数なら将来価値を現在決定へ戻す役割、FIFOなら旅行時間情報モデルの整合性を保証する役割、CFAなら目的関数を調整して将来に良い行動を誘導する役割である。\end{enumerate}
\end{minipage}
\end{adjustbox}
\end{minipage}


\clearpage
\SlideHeader{Lecture 11 - Platform-based Food Delivery}{019}{Modeling (4)}
\begin{minipage}[t]{0.405\textwidth}
\SlideImage{19}{../Materials/2026/3DM_11_Food-Delivery_SS26.pdf}
\begin{adjustbox}{max totalsize={\textwidth}{0.43\textheight},center}
\begin{minipage}{\textwidth}\scriptsize
\NoteSection{日本語訳}\begin{itemize}
\item Modeling (4)（この用語は講義スライド上の専門語として扱う）
\item ▪ Exploration by trying a new parameter set is costly in terms of time and profit (loss)（この用語は講義スライド上の専門語として扱う）
\item ▪ The number of parameter sets to try is very limited → Bayesian optimization applies（この用語は講義スライド上の専門語として扱う）
\item ▪ At t=3 three sets have been tested（この用語は講義スライド上の専門語として扱う）
\item ▪ For these sets we‘ve got confidence（この用語は講義スライド上の専門語として扱う）
\item ▪ An acquisition function suggests new trial（この用語は講義スライド上の専門語として扱う）
\item ▪ At t=4 the confidence has improved（この用語は講義スライド上の専門語として扱う）
\item because of the new observation（この用語は講義スライド上の専門語として扱う）
\item ▪ A modified acquisition function suggests（この用語は講義スライド上の専門語として扱う）
\end{itemize}\NoteSection{試験対策ポイント}\begin{itemize}
\item Platform-based food deliveryは、需要・供給・推薦・配送の複数意思決定が連鎖する問題である。
\item food recommendationとdelivery optimizationを分けつつ、相互作用を説明する。
\end{itemize}
\end{minipage}
\end{adjustbox}
\end{minipage}\hfill
\begin{minipage}[t]{0.58\textwidth}
\begin{adjustbox}{max totalsize={\textwidth}{0.89\textheight},center}
\begin{minipage}{\textwidth}\scriptsize
\NoteSection{解説}このページでは「Modeling (4)」を理解する。スライド上の表現は短いが、試験では定義、具体例、モデル上の意味、手法の限界まで説明できる必要がある。\par
Platform-based food delivery は、顧客、レストラン、配達員、プラットフォームが相互作用する多面的なサービスである。顧客は料理を選び、レストランは注文を調理し、配達員は複数注文を運び、プラットフォームは推薦、価格、割当、ルーティングを決める。\par
この問題では、単に最短ルートを作るだけでは不十分である。推薦システムが特定店舗に需要を集中させると、調理待ちや配達遅延が増える可能性がある。逆に、配送能力を考慮した推薦を行えば、需要を空いている店舗や配送しやすい地域へ誘導できる。\par
food recommendation では、content-based recommendation, collaborative filtering, ontology graph, diversificationなどが扱われる。顧客の好みだけでなく、料理の種類、価格、距離、調理時間、配達可能性などが推薦品質に影響する。\par
配送最適化では、注文受諾、配達員割当、ピックアップ順、ドロップオフ順、複数注文のバッチングが問題になる。時間窓、料理の品質劣化、顧客待ち時間、配達員の公平性なども考慮する必要がある。\par
試験では、食品配送をMDP/ADPとして設計する場合、状態に未処理注文、配達員位置、レストラン調理状態、需要予測を入れ、決定に推薦・注文受諾・割当・ルーティングを入れる。報酬には収益、遅延ペナルティ、顧客満足度、配達員稼働を含めるとよい。\par\NoteSection{確認問題と解答}\begin{enumerate}\item \textbf{L11-S019: 「Modeling (4)」の概念を、定義・具体例・モデル上の意味の順に説明せよ。}\\定義としては、Platform-based food deliveryは、需要・供給・推薦・配送の複数意思決定が連鎖する問題である。。具体例では、配送・在庫・フードデリバリーのような現実問題を用い、状態、決定、外生情報、報酬、または情報モデル/意思決定モデルのどこに対応するかを明示する。モデル上の意味として、この概念が目的関数、制約、状態表現、将来価値、または方策選択にどのような影響を与えるかを書く。試験では、用語だけを列挙せず、入力データから意思決定、さらに現実の実装へつながる流れを文章で示すことが重要である。\item \textbf{L11-S019: このスライドの考え方を、講義全体のDDDMプロセスの中に位置づけよ。}\\DDDMプロセスでは、現実世界のデータを集約して情報モデルを作り、意思決定モデルまたは方策で行動を選び、実装後に結果を観測して次の状態へ進む。このスライドの概念は、そのどこか一部だけではなく、データ表現、意思決定、将来不確実性、計算可能性のいずれかに関わる。答案では、まずこの概念がどの段階に属するかを述べ、次に他の段階へどう影響するかを書く。例えば価値関数なら将来価値を現在決定へ戻す役割、FIFOなら旅行時間情報モデルの整合性を保証する役割、CFAなら目的関数を調整して将来に良い行動を誘導する役割である。\end{enumerate}
\end{minipage}
\end{adjustbox}
\end{minipage}


\clearpage
\SlideHeader{Lecture 11 - Platform-based Food Delivery}{020}{Prediction and optimization pipeline}
\begin{minipage}[t]{0.405\textwidth}
\SlideImage{20}{../Materials/2026/3DM_11_Food-Delivery_SS26.pdf}
\begin{adjustbox}{max totalsize={\textwidth}{0.43\textheight},center}
\begin{minipage}{\textwidth}\scriptsize
\NoteSection{日本語訳}\begin{itemize}
\item Prediction and optimization pipeline（この用語は講義スライド上の専門語として扱う）
\item ▪ Data analysis predicts demand (fulfilment) in the next period（この用語は講義スライド上の専門語として扱う）
\item ▪ Multi-objective optimization identifies appropriate restaurants/meals（この用語は講義スライド上の専門語として扱う）
\item https://www.uber.com/newsroom/arriving-now-the-new-uber-eats/（この用語は講義スライド上の専門語として扱う）
\item ▪ Eater‘s eyeball triggers order that 最大化するs conversion rate
\item https://www.uber.com/en-DE/blog/uber-eats-recommending-marketplace/（この用語は講義スライド上の専門語として扱う）
\end{itemize}\NoteSection{試験対策ポイント}\begin{itemize}
\item Platform-based food deliveryは、需要・供給・推薦・配送の複数意思決定が連鎖する問題である。
\item food recommendationとdelivery optimizationを分けつつ、相互作用を説明する。
\end{itemize}
\end{minipage}
\end{adjustbox}
\end{minipage}\hfill
\begin{minipage}[t]{0.58\textwidth}
\begin{adjustbox}{max totalsize={\textwidth}{0.89\textheight},center}
\begin{minipage}{\textwidth}\scriptsize
\NoteSection{解説}このページでは「Prediction and optimization pipeline」を理解する。スライド上の表現は短いが、試験では定義、具体例、モデル上の意味、手法の限界まで説明できる必要がある。\par
Platform-based food delivery は、顧客、レストラン、配達員、プラットフォームが相互作用する多面的なサービスである。顧客は料理を選び、レストランは注文を調理し、配達員は複数注文を運び、プラットフォームは推薦、価格、割当、ルーティングを決める。\par
この問題では、単に最短ルートを作るだけでは不十分である。推薦システムが特定店舗に需要を集中させると、調理待ちや配達遅延が増える可能性がある。逆に、配送能力を考慮した推薦を行えば、需要を空いている店舗や配送しやすい地域へ誘導できる。\par
food recommendation では、content-based recommendation, collaborative filtering, ontology graph, diversificationなどが扱われる。顧客の好みだけでなく、料理の種類、価格、距離、調理時間、配達可能性などが推薦品質に影響する。\par
配送最適化では、注文受諾、配達員割当、ピックアップ順、ドロップオフ順、複数注文のバッチングが問題になる。時間窓、料理の品質劣化、顧客待ち時間、配達員の公平性なども考慮する必要がある。\par
試験では、食品配送をMDP/ADPとして設計する場合、状態に未処理注文、配達員位置、レストラン調理状態、需要予測を入れ、決定に推薦・注文受諾・割当・ルーティングを入れる。報酬には収益、遅延ペナルティ、顧客満足度、配達員稼働を含めるとよい。\par\NoteSection{確認問題と解答}\begin{enumerate}\item \textbf{L11-S020: 「Prediction and optimization pipeline」の概念を、定義・具体例・モデル上の意味の順に説明せよ。}\\定義としては、Platform-based food deliveryは、需要・供給・推薦・配送の複数意思決定が連鎖する問題である。。具体例では、配送・在庫・フードデリバリーのような現実問題を用い、状態、決定、外生情報、報酬、または情報モデル/意思決定モデルのどこに対応するかを明示する。モデル上の意味として、この概念が目的関数、制約、状態表現、将来価値、または方策選択にどのような影響を与えるかを書く。試験では、用語だけを列挙せず、入力データから意思決定、さらに現実の実装へつながる流れを文章で示すことが重要である。\item \textbf{L11-S020: このスライドの考え方を、講義全体のDDDMプロセスの中に位置づけよ。}\\DDDMプロセスでは、現実世界のデータを集約して情報モデルを作り、意思決定モデルまたは方策で行動を選び、実装後に結果を観測して次の状態へ進む。このスライドの概念は、そのどこか一部だけではなく、データ表現、意思決定、将来不確実性、計算可能性のいずれかに関わる。答案では、まずこの概念がどの段階に属するかを述べ、次に他の段階へどう影響するかを書く。例えば価値関数なら将来価値を現在決定へ戻す役割、FIFOなら旅行時間情報モデルの整合性を保証する役割、CFAなら目的関数を調整して将来に良い行動を誘導する役割である。\end{enumerate}
\end{minipage}
\end{adjustbox}
\end{minipage}


\clearpage
\SlideHeader{Lecture 11 - Platform-based Food Delivery}{021}{Content}
\begin{minipage}[t]{0.405\textwidth}
\SlideImage{21}{../Materials/2026/3DM_11_Food-Delivery_SS26.pdf}
\begin{adjustbox}{max totalsize={\textwidth}{0.43\textheight},center}
\begin{minipage}{\textwidth}\scriptsize
\NoteSection{日本語訳}\begin{itemize}
\item 内容
\item ▪ Platform-based Food Delivery（この用語は講義スライド上の専門語として扱う）
\item ▪ The Eater‘s Perspective（この用語は講義スライド上の専門語として扱う）
\item ▪ Uber Eat’s Perspective（この用語は講義スライド上の専門語として扱う）
\item ▪ The Delivery Perspective（この用語は講義スライド上の専門語として扱う）
\item ▪ The Restaurant‘s Perspective（この用語は講義スライド上の専門語として扱う）
\end{itemize}
\end{minipage}
\end{adjustbox}
\end{minipage}\hfill
\begin{minipage}[t]{0.58\textwidth}
\begin{adjustbox}{max totalsize={\textwidth}{0.89\textheight},center}
\begin{minipage}{\textwidth}\scriptsize
\NoteSection{注記}このページは表紙・目次・事務情報・導入画像が中心であるため、無理に確認問題や過去問を付けない。
\end{minipage}
\end{adjustbox}
\end{minipage}


\clearpage
\SlideHeader{Lecture 11 - Platform-based Food Delivery}{022}{Interaction of partners}
\begin{minipage}[t]{0.405\textwidth}
\SlideImage{22}{../Materials/2026/3DM_11_Food-Delivery_SS26.pdf}
\begin{adjustbox}{max totalsize={\textwidth}{0.43\textheight},center}
\begin{minipage}{\textwidth}\scriptsize
\NoteSection{日本語訳}\begin{itemize}
\item Interaction of partners（この用語は講義スライド上の専門語として扱う）
\item https://www.infoq.com/fr/articles/uber-eats-time-predictions/（この用語は講義スライド上の専門語として扱う）
\end{itemize}\NoteSection{試験対策ポイント}\begin{itemize}
\item Platform-based food deliveryは、需要・供給・推薦・配送の複数意思決定が連鎖する問題である。
\item food recommendationとdelivery optimizationを分けつつ、相互作用を説明する。
\end{itemize}
\end{minipage}
\end{adjustbox}
\end{minipage}\hfill
\begin{minipage}[t]{0.58\textwidth}
\begin{adjustbox}{max totalsize={\textwidth}{0.89\textheight},center}
\begin{minipage}{\textwidth}\scriptsize
\NoteSection{解説}このページでは「Interaction of partners」を理解する。スライド上の表現は短いが、試験では定義、具体例、モデル上の意味、手法の限界まで説明できる必要がある。\par
Platform-based food delivery は、顧客、レストラン、配達員、プラットフォームが相互作用する多面的なサービスである。顧客は料理を選び、レストランは注文を調理し、配達員は複数注文を運び、プラットフォームは推薦、価格、割当、ルーティングを決める。\par
この問題では、単に最短ルートを作るだけでは不十分である。推薦システムが特定店舗に需要を集中させると、調理待ちや配達遅延が増える可能性がある。逆に、配送能力を考慮した推薦を行えば、需要を空いている店舗や配送しやすい地域へ誘導できる。\par
food recommendation では、content-based recommendation, collaborative filtering, ontology graph, diversificationなどが扱われる。顧客の好みだけでなく、料理の種類、価格、距離、調理時間、配達可能性などが推薦品質に影響する。\par
配送最適化では、注文受諾、配達員割当、ピックアップ順、ドロップオフ順、複数注文のバッチングが問題になる。時間窓、料理の品質劣化、顧客待ち時間、配達員の公平性なども考慮する必要がある。\par
試験では、食品配送をMDP/ADPとして設計する場合、状態に未処理注文、配達員位置、レストラン調理状態、需要予測を入れ、決定に推薦・注文受諾・割当・ルーティングを入れる。報酬には収益、遅延ペナルティ、顧客満足度、配達員稼働を含めるとよい。\par\NoteSection{確認問題と解答}\begin{enumerate}\item \textbf{L11-S022: 「Interaction of partners」の概念を、定義・具体例・モデル上の意味の順に説明せよ。}\\定義としては、Platform-based food deliveryは、需要・供給・推薦・配送の複数意思決定が連鎖する問題である。。具体例では、配送・在庫・フードデリバリーのような現実問題を用い、状態、決定、外生情報、報酬、または情報モデル/意思決定モデルのどこに対応するかを明示する。モデル上の意味として、この概念が目的関数、制約、状態表現、将来価値、または方策選択にどのような影響を与えるかを書く。試験では、用語だけを列挙せず、入力データから意思決定、さらに現実の実装へつながる流れを文章で示すことが重要である。\item \textbf{L11-S022: このスライドの考え方を、講義全体のDDDMプロセスの中に位置づけよ。}\\DDDMプロセスでは、現実世界のデータを集約して情報モデルを作り、意思決定モデルまたは方策で行動を選び、実装後に結果を観測して次の状態へ進む。このスライドの概念は、そのどこか一部だけではなく、データ表現、意思決定、将来不確実性、計算可能性のいずれかに関わる。答案では、まずこの概念がどの段階に属するかを述べ、次に他の段階へどう影響するかを書く。例えば価値関数なら将来価値を現在決定へ戻す役割、FIFOなら旅行時間情報モデルの整合性を保証する役割、CFAなら目的関数を調整して将来に良い行動を誘導する役割である。\end{enumerate}
\end{minipage}
\end{adjustbox}
\end{minipage}


\clearpage
\SlideHeader{Lecture 11 - Platform-based Food Delivery}{023}{Dispatching a driver}
\begin{minipage}[t]{0.405\textwidth}
\SlideImage{23}{../Materials/2026/3DM_11_Food-Delivery_SS26.pdf}
\begin{adjustbox}{max totalsize={\textwidth}{0.43\textheight},center}
\begin{minipage}{\textwidth}\scriptsize
\NoteSection{日本語訳}\begin{itemize}
\item Dispatching a driver（この用語は講義スライド上の専門語として扱う）
\item ▪ Given a (30min) meal preparation time, dispatch of drivers starts 10min before ETA（この用語は講義スライド上の専門語として扱う）
\item ▪ Several dispatch attempts until successful trip acceptance by one of the driver pool（この用語は講義スライド上の専門語として扱う）
\item ▪ Drivers may be ranked and (pre-)selected prior to dispatch attempts（この用語は講義スライド上の専門語として扱う）
\item https://www.infoq.com/fr/articles/uber-eats-time-predictions/（この用語は講義スライド上の専門語として扱う）
\end{itemize}\NoteSection{試験対策ポイント}\begin{itemize}
\item Platform-based food deliveryは、需要・供給・推薦・配送の複数意思決定が連鎖する問題である。
\item food recommendationとdelivery optimizationを分けつつ、相互作用を説明する。
\end{itemize}
\end{minipage}
\end{adjustbox}
\end{minipage}\hfill
\begin{minipage}[t]{0.58\textwidth}
\begin{adjustbox}{max totalsize={\textwidth}{0.89\textheight},center}
\begin{minipage}{\textwidth}\scriptsize
\NoteSection{解説}このページでは「Dispatching a driver」を理解する。スライド上の表現は短いが、試験では定義、具体例、モデル上の意味、手法の限界まで説明できる必要がある。\par
Platform-based food delivery は、顧客、レストラン、配達員、プラットフォームが相互作用する多面的なサービスである。顧客は料理を選び、レストランは注文を調理し、配達員は複数注文を運び、プラットフォームは推薦、価格、割当、ルーティングを決める。\par
この問題では、単に最短ルートを作るだけでは不十分である。推薦システムが特定店舗に需要を集中させると、調理待ちや配達遅延が増える可能性がある。逆に、配送能力を考慮した推薦を行えば、需要を空いている店舗や配送しやすい地域へ誘導できる。\par
food recommendation では、content-based recommendation, collaborative filtering, ontology graph, diversificationなどが扱われる。顧客の好みだけでなく、料理の種類、価格、距離、調理時間、配達可能性などが推薦品質に影響する。\par
配送最適化では、注文受諾、配達員割当、ピックアップ順、ドロップオフ順、複数注文のバッチングが問題になる。時間窓、料理の品質劣化、顧客待ち時間、配達員の公平性なども考慮する必要がある。\par
試験では、食品配送をMDP/ADPとして設計する場合、状態に未処理注文、配達員位置、レストラン調理状態、需要予測を入れ、決定に推薦・注文受諾・割当・ルーティングを入れる。報酬には収益、遅延ペナルティ、顧客満足度、配達員稼働を含めるとよい。\par\NoteSection{確認問題と解答}\begin{enumerate}\item \textbf{L11-S023: 「Dispatching a driver」の概念を、定義・具体例・モデル上の意味の順に説明せよ。}\\定義としては、Platform-based food deliveryは、需要・供給・推薦・配送の複数意思決定が連鎖する問題である。。具体例では、配送・在庫・フードデリバリーのような現実問題を用い、状態、決定、外生情報、報酬、または情報モデル/意思決定モデルのどこに対応するかを明示する。モデル上の意味として、この概念が目的関数、制約、状態表現、将来価値、または方策選択にどのような影響を与えるかを書く。試験では、用語だけを列挙せず、入力データから意思決定、さらに現実の実装へつながる流れを文章で示すことが重要である。\item \textbf{L11-S023: このスライドの考え方を、講義全体のDDDMプロセスの中に位置づけよ。}\\DDDMプロセスでは、現実世界のデータを集約して情報モデルを作り、意思決定モデルまたは方策で行動を選び、実装後に結果を観測して次の状態へ進む。このスライドの概念は、そのどこか一部だけではなく、データ表現、意思決定、将来不確実性、計算可能性のいずれかに関わる。答案では、まずこの概念がどの段階に属するかを述べ、次に他の段階へどう影響するかを書く。例えば価値関数なら将来価値を現在決定へ戻す役割、FIFOなら旅行時間情報モデルの整合性を保証する役割、CFAなら目的関数を調整して将来に良い行動を誘導する役割である。\end{enumerate}
\end{minipage}
\end{adjustbox}
\end{minipage}


\clearpage
\SlideHeader{Lecture 11 - Platform-based Food Delivery}{024}{Driver selection}
\begin{minipage}[t]{0.405\textwidth}
\SlideImage{24}{../Materials/2026/3DM_11_Food-Delivery_SS26.pdf}
\begin{adjustbox}{max totalsize={\textwidth}{0.43\textheight},center}
\begin{minipage}{\textwidth}\scriptsize
\NoteSection{日本語訳}\begin{itemize}
\item Driver selection（この用語は講義スライド上の専門語として扱う）
\item https://www.infoq.com/fr/articles/uber-eats-time-predictions/（この用語は講義スライド上の専門語として扱う）
\item ▪ Distance to the restaurant（この用語は講義スライド上の専門語として扱う）
\item ▪ Second order in vicinity（この用語は講義スライド上の専門語として扱う）
\item ▪ Reliability of the driver（この用語は講義スライド上の専門語として扱う）
\item ▪ Working time restriction（この用語は講義スライド上の専門語として扱う）
\item ▪ Number of recent deliveries（この用語は講義スライド上の専門語として扱う）
\item ▪ Weekly payout so far（この用語は講義スライド上の専門語として扱う）
\item ▪ Batched opportunities:（この用語は講義スライド上の専門語として扱う）
\end{itemize}\NoteSection{試験対策ポイント}\begin{itemize}
\item Platform-based food deliveryは、需要・供給・推薦・配送の複数意思決定が連鎖する問題である。
\item food recommendationとdelivery optimizationを分けつつ、相互作用を説明する。
\end{itemize}
\end{minipage}
\end{adjustbox}
\end{minipage}\hfill
\begin{minipage}[t]{0.58\textwidth}
\begin{adjustbox}{max totalsize={\textwidth}{0.89\textheight},center}
\begin{minipage}{\textwidth}\scriptsize
\NoteSection{解説}このページでは「Driver selection」を理解する。スライド上の表現は短いが、試験では定義、具体例、モデル上の意味、手法の限界まで説明できる必要がある。\par
Platform-based food delivery は、顧客、レストラン、配達員、プラットフォームが相互作用する多面的なサービスである。顧客は料理を選び、レストランは注文を調理し、配達員は複数注文を運び、プラットフォームは推薦、価格、割当、ルーティングを決める。\par
この問題では、単に最短ルートを作るだけでは不十分である。推薦システムが特定店舗に需要を集中させると、調理待ちや配達遅延が増える可能性がある。逆に、配送能力を考慮した推薦を行えば、需要を空いている店舗や配送しやすい地域へ誘導できる。\par
food recommendation では、content-based recommendation, collaborative filtering, ontology graph, diversificationなどが扱われる。顧客の好みだけでなく、料理の種類、価格、距離、調理時間、配達可能性などが推薦品質に影響する。\par
配送最適化では、注文受諾、配達員割当、ピックアップ順、ドロップオフ順、複数注文のバッチングが問題になる。時間窓、料理の品質劣化、顧客待ち時間、配達員の公平性なども考慮する必要がある。\par
試験では、食品配送をMDP/ADPとして設計する場合、状態に未処理注文、配達員位置、レストラン調理状態、需要予測を入れ、決定に推薦・注文受諾・割当・ルーティングを入れる。報酬には収益、遅延ペナルティ、顧客満足度、配達員稼働を含めるとよい。\par\NoteSection{確認問題と解答}\begin{enumerate}\item \textbf{L11-S024: 「Driver selection」の概念を、定義・具体例・モデル上の意味の順に説明せよ。}\\定義としては、Platform-based food deliveryは、需要・供給・推薦・配送の複数意思決定が連鎖する問題である。。具体例では、配送・在庫・フードデリバリーのような現実問題を用い、状態、決定、外生情報、報酬、または情報モデル/意思決定モデルのどこに対応するかを明示する。モデル上の意味として、この概念が目的関数、制約、状態表現、将来価値、または方策選択にどのような影響を与えるかを書く。試験では、用語だけを列挙せず、入力データから意思決定、さらに現実の実装へつながる流れを文章で示すことが重要である。\item \textbf{L11-S024: このスライドの考え方を、講義全体のDDDMプロセスの中に位置づけよ。}\\DDDMプロセスでは、現実世界のデータを集約して情報モデルを作り、意思決定モデルまたは方策で行動を選び、実装後に結果を観測して次の状態へ進む。このスライドの概念は、そのどこか一部だけではなく、データ表現、意思決定、将来不確実性、計算可能性のいずれかに関わる。答案では、まずこの概念がどの段階に属するかを述べ、次に他の段階へどう影響するかを書く。例えば価値関数なら将来価値を現在決定へ戻す役割、FIFOなら旅行時間情報モデルの整合性を保証する役割、CFAなら目的関数を調整して将来に良い行動を誘導する役割である。\end{enumerate}
\end{minipage}
\end{adjustbox}
\end{minipage}


\clearpage
\SlideHeader{Lecture 11 - Platform-based Food Delivery}{025}{State of order is communicated to the eater}
\begin{minipage}[t]{0.405\textwidth}
\SlideImage{25}{../Materials/2026/3DM_11_Food-Delivery_SS26.pdf}
\begin{adjustbox}{max totalsize={\textwidth}{0.43\textheight},center}
\begin{minipage}{\textwidth}\scriptsize
\NoteSection{日本語訳}\begin{itemize}
\item 状態 of order is communicated to the eater
\item ▪ The driver is monitored（この用語は講義スライド上の専門語として扱う）
\item ▪ The position of the driver is（この用語は講義スライド上の専門語として扱う）
\item communicated to the eater（この用語は講義スライド上の専門語として扱う）
\item ▪ The predicted time of delivery is（この用語は講義スライド上の専門語として扱う）
\item updated（この用語は講義スライド上の専門語として扱う）
\item ▪ Drivers are ranked by the platform（この用語は講義スライド上の専門語として扱う）
\item ▪ By the speed of delivery（この用語は講義スライド上の専門語として扱う）
\item ▪ By the reliability of activities（この用語は講義スライド上の専門語として扱う）
\end{itemize}\NoteSection{試験対策ポイント}\begin{itemize}
\item Platform-based food deliveryは、需要・供給・推薦・配送の複数意思決定が連鎖する問題である。
\item food recommendationとdelivery optimizationを分けつつ、相互作用を説明する。
\end{itemize}
\end{minipage}
\end{adjustbox}
\end{minipage}\hfill
\begin{minipage}[t]{0.58\textwidth}
\begin{adjustbox}{max totalsize={\textwidth}{0.89\textheight},center}
\begin{minipage}{\textwidth}\scriptsize
\NoteSection{解説}このページでは「State of order is communicated to the eater」を理解する。スライド上の表現は短いが、試験では定義、具体例、モデル上の意味、手法の限界まで説明できる必要がある。\par
Platform-based food delivery は、顧客、レストラン、配達員、プラットフォームが相互作用する多面的なサービスである。顧客は料理を選び、レストランは注文を調理し、配達員は複数注文を運び、プラットフォームは推薦、価格、割当、ルーティングを決める。\par
この問題では、単に最短ルートを作るだけでは不十分である。推薦システムが特定店舗に需要を集中させると、調理待ちや配達遅延が増える可能性がある。逆に、配送能力を考慮した推薦を行えば、需要を空いている店舗や配送しやすい地域へ誘導できる。\par
food recommendation では、content-based recommendation, collaborative filtering, ontology graph, diversificationなどが扱われる。顧客の好みだけでなく、料理の種類、価格、距離、調理時間、配達可能性などが推薦品質に影響する。\par
配送最適化では、注文受諾、配達員割当、ピックアップ順、ドロップオフ順、複数注文のバッチングが問題になる。時間窓、料理の品質劣化、顧客待ち時間、配達員の公平性なども考慮する必要がある。\par
試験では、食品配送をMDP/ADPとして設計する場合、状態に未処理注文、配達員位置、レストラン調理状態、需要予測を入れ、決定に推薦・注文受諾・割当・ルーティングを入れる。報酬には収益、遅延ペナルティ、顧客満足度、配達員稼働を含めるとよい。\par\NoteSection{確認問題と解答}\begin{enumerate}\item \textbf{L11-S025: 「State of order is communicated to the eater」の概念を、定義・具体例・モデル上の意味の順に説明せよ。}\\定義としては、Platform-based food deliveryは、需要・供給・推薦・配送の複数意思決定が連鎖する問題である。。具体例では、配送・在庫・フードデリバリーのような現実問題を用い、状態、決定、外生情報、報酬、または情報モデル/意思決定モデルのどこに対応するかを明示する。モデル上の意味として、この概念が目的関数、制約、状態表現、将来価値、または方策選択にどのような影響を与えるかを書く。試験では、用語だけを列挙せず、入力データから意思決定、さらに現実の実装へつながる流れを文章で示すことが重要である。\item \textbf{L11-S025: このスライドの考え方を、講義全体のDDDMプロセスの中に位置づけよ。}\\DDDMプロセスでは、現実世界のデータを集約して情報モデルを作り、意思決定モデルまたは方策で行動を選び、実装後に結果を観測して次の状態へ進む。このスライドの概念は、そのどこか一部だけではなく、データ表現、意思決定、将来不確実性、計算可能性のいずれかに関わる。答案では、まずこの概念がどの段階に属するかを述べ、次に他の段階へどう影響するかを書く。例えば価値関数なら将来価値を現在決定へ戻す役割、FIFOなら旅行時間情報モデルの整合性を保証する役割、CFAなら目的関数を調整して将来に良い行動を誘導する役割である。\end{enumerate}
\end{minipage}
\end{adjustbox}
\end{minipage}


\clearpage
\SlideHeader{Lecture 11 - Platform-based Food Delivery}{026}{Content}
\begin{minipage}[t]{0.405\textwidth}
\SlideImage{26}{../Materials/2026/3DM_11_Food-Delivery_SS26.pdf}
\begin{adjustbox}{max totalsize={\textwidth}{0.43\textheight},center}
\begin{minipage}{\textwidth}\scriptsize
\NoteSection{日本語訳}\begin{itemize}
\item 内容
\item ▪ Platform-based Food Delivery（この用語は講義スライド上の専門語として扱う）
\item ▪ The Eater‘s Perspective（この用語は講義スライド上の専門語として扱う）
\item ▪ Uber Eat’s Perspective（この用語は講義スライド上の専門語として扱う）
\item ▪ The Delivery Perspective（この用語は講義スライド上の専門語として扱う）
\item ▪ The Restaurant‘s Perspective（この用語は講義スライド上の専門語として扱う）
\end{itemize}
\end{minipage}
\end{adjustbox}
\end{minipage}\hfill
\begin{minipage}[t]{0.58\textwidth}
\begin{adjustbox}{max totalsize={\textwidth}{0.89\textheight},center}
\begin{minipage}{\textwidth}\scriptsize
\NoteSection{注記}このページは表紙・目次・事務情報・導入画像が中心であるため、無理に確認問題や過去問を付けない。
\end{minipage}
\end{adjustbox}
\end{minipage}


\clearpage
\SlideHeader{Lecture 11 - Platform-based Food Delivery}{027}{Interaction of partner}
\begin{minipage}[t]{0.405\textwidth}
\SlideImage{27}{../Materials/2026/3DM_11_Food-Delivery_SS26.pdf}
\begin{adjustbox}{max totalsize={\textwidth}{0.43\textheight},center}
\begin{minipage}{\textwidth}\scriptsize
\NoteSection{日本語訳}\begin{itemize}
\item 参加者間の相互作用
\item https://www.infoq.com/fr/articles/uber-eats-time-predictions/（この用語は講義スライド上の専門語として扱う）
\end{itemize}\NoteSection{試験対策ポイント}\begin{itemize}
\item Platform-based food deliveryは、需要・供給・推薦・配送の複数意思決定が連鎖する問題である。
\item food recommendationとdelivery optimizationを分けつつ、相互作用を説明する。
\end{itemize}
\end{minipage}
\end{adjustbox}
\end{minipage}\hfill
\begin{minipage}[t]{0.58\textwidth}
\begin{adjustbox}{max totalsize={\textwidth}{0.89\textheight},center}
\begin{minipage}{\textwidth}\scriptsize
\NoteSection{解説}このページでは「Interaction of partner」を理解する。スライド上の表現は短いが、試験では定義、具体例、モデル上の意味、手法の限界まで説明できる必要がある。\par
Platform-based food delivery は、顧客、レストラン、配達員、プラットフォームが相互作用する多面的なサービスである。顧客は料理を選び、レストランは注文を調理し、配達員は複数注文を運び、プラットフォームは推薦、価格、割当、ルーティングを決める。\par
この問題では、単に最短ルートを作るだけでは不十分である。推薦システムが特定店舗に需要を集中させると、調理待ちや配達遅延が増える可能性がある。逆に、配送能力を考慮した推薦を行えば、需要を空いている店舗や配送しやすい地域へ誘導できる。\par
food recommendation では、content-based recommendation, collaborative filtering, ontology graph, diversificationなどが扱われる。顧客の好みだけでなく、料理の種類、価格、距離、調理時間、配達可能性などが推薦品質に影響する。\par
配送最適化では、注文受諾、配達員割当、ピックアップ順、ドロップオフ順、複数注文のバッチングが問題になる。時間窓、料理の品質劣化、顧客待ち時間、配達員の公平性なども考慮する必要がある。\par
試験では、食品配送をMDP/ADPとして設計する場合、状態に未処理注文、配達員位置、レストラン調理状態、需要予測を入れ、決定に推薦・注文受諾・割当・ルーティングを入れる。報酬には収益、遅延ペナルティ、顧客満足度、配達員稼働を含めるとよい。\par\NoteSection{確認問題と解答}\begin{enumerate}\item \textbf{L11-S027: 「Interaction of partner」の概念を、定義・具体例・モデル上の意味の順に説明せよ。}\\定義としては、Platform-based food deliveryは、需要・供給・推薦・配送の複数意思決定が連鎖する問題である。。具体例では、配送・在庫・フードデリバリーのような現実問題を用い、状態、決定、外生情報、報酬、または情報モデル/意思決定モデルのどこに対応するかを明示する。モデル上の意味として、この概念が目的関数、制約、状態表現、将来価値、または方策選択にどのような影響を与えるかを書く。試験では、用語だけを列挙せず、入力データから意思決定、さらに現実の実装へつながる流れを文章で示すことが重要である。\item \textbf{L11-S027: このスライドの考え方を、講義全体のDDDMプロセスの中に位置づけよ。}\\DDDMプロセスでは、現実世界のデータを集約して情報モデルを作り、意思決定モデルまたは方策で行動を選び、実装後に結果を観測して次の状態へ進む。このスライドの概念は、そのどこか一部だけではなく、データ表現、意思決定、将来不確実性、計算可能性のいずれかに関わる。答案では、まずこの概念がどの段階に属するかを述べ、次に他の段階へどう影響するかを書く。例えば価値関数なら将来価値を現在決定へ戻す役割、FIFOなら旅行時間情報モデルの整合性を保証する役割、CFAなら目的関数を調整して将来に良い行動を誘導する役割である。\end{enumerate}
\end{minipage}
\end{adjustbox}
\end{minipage}


\clearpage
\SlideHeader{Lecture 11 - Platform-based Food Delivery}{028}{Food preparation at the restaurant}
\begin{minipage}[t]{0.405\textwidth}
\SlideImage{28}{../Materials/2026/3DM_11_Food-Delivery_SS26.pdf}
\begin{adjustbox}{max totalsize={\textwidth}{0.43\textheight},center}
\begin{minipage}{\textwidth}\scriptsize
\NoteSection{日本語訳}\begin{itemize}
\item Food preparation at the restaurant（この用語は講義スライド上の専門語として扱う）
\item ▪ After order, speedy meal preparation and delivery is of utmost importance（この用語は講義スライド上の専門語として扱う）
\item ▪ A crucial point of time is the food handover to the driver that cannot be observed（この用語は講義スライド上の専門語として扱う）
\item https://www.uber.com/en-DE/blog/uber-eats-recommending-（この用語は講義スライド上の専門語として扱う）
\item marketplace/（この用語は講義スライド上の専門語として扱う）
\end{itemize}\NoteSection{試験対策ポイント}\begin{itemize}
\item Platform-based food deliveryは、需要・供給・推薦・配送の複数意思決定が連鎖する問題である。
\item food recommendationとdelivery optimizationを分けつつ、相互作用を説明する。
\end{itemize}
\end{minipage}
\end{adjustbox}
\end{minipage}\hfill
\begin{minipage}[t]{0.58\textwidth}
\begin{adjustbox}{max totalsize={\textwidth}{0.89\textheight},center}
\begin{minipage}{\textwidth}\scriptsize
\NoteSection{解説}このページでは「Food preparation at the restaurant」を理解する。スライド上の表現は短いが、試験では定義、具体例、モデル上の意味、手法の限界まで説明できる必要がある。\par
Platform-based food delivery は、顧客、レストラン、配達員、プラットフォームが相互作用する多面的なサービスである。顧客は料理を選び、レストランは注文を調理し、配達員は複数注文を運び、プラットフォームは推薦、価格、割当、ルーティングを決める。\par
この問題では、単に最短ルートを作るだけでは不十分である。推薦システムが特定店舗に需要を集中させると、調理待ちや配達遅延が増える可能性がある。逆に、配送能力を考慮した推薦を行えば、需要を空いている店舗や配送しやすい地域へ誘導できる。\par
food recommendation では、content-based recommendation, collaborative filtering, ontology graph, diversificationなどが扱われる。顧客の好みだけでなく、料理の種類、価格、距離、調理時間、配達可能性などが推薦品質に影響する。\par
配送最適化では、注文受諾、配達員割当、ピックアップ順、ドロップオフ順、複数注文のバッチングが問題になる。時間窓、料理の品質劣化、顧客待ち時間、配達員の公平性なども考慮する必要がある。\par
試験では、食品配送をMDP/ADPとして設計する場合、状態に未処理注文、配達員位置、レストラン調理状態、需要予測を入れ、決定に推薦・注文受諾・割当・ルーティングを入れる。報酬には収益、遅延ペナルティ、顧客満足度、配達員稼働を含めるとよい。\par\NoteSection{確認問題と解答}\begin{enumerate}\item \textbf{L11-S028: 「Food preparation at the restaurant」の概念を、定義・具体例・モデル上の意味の順に説明せよ。}\\定義としては、Platform-based food deliveryは、需要・供給・推薦・配送の複数意思決定が連鎖する問題である。。具体例では、配送・在庫・フードデリバリーのような現実問題を用い、状態、決定、外生情報、報酬、または情報モデル/意思決定モデルのどこに対応するかを明示する。モデル上の意味として、この概念が目的関数、制約、状態表現、将来価値、または方策選択にどのような影響を与えるかを書く。試験では、用語だけを列挙せず、入力データから意思決定、さらに現実の実装へつながる流れを文章で示すことが重要である。\item \textbf{L11-S028: このスライドの考え方を、講義全体のDDDMプロセスの中に位置づけよ。}\\DDDMプロセスでは、現実世界のデータを集約して情報モデルを作り、意思決定モデルまたは方策で行動を選び、実装後に結果を観測して次の状態へ進む。このスライドの概念は、そのどこか一部だけではなく、データ表現、意思決定、将来不確実性、計算可能性のいずれかに関わる。答案では、まずこの概念がどの段階に属するかを述べ、次に他の段階へどう影響するかを書く。例えば価値関数なら将来価値を現在決定へ戻す役割、FIFOなら旅行時間情報モデルの整合性を保証する役割、CFAなら目的関数を調整して将来に良い行動を誘導する役割である。\end{enumerate}
\end{minipage}
\end{adjustbox}
\end{minipage}


\clearpage
\SlideHeader{Lecture 11 - Platform-based Food Delivery}{029}{Online and offline prediction of food handover time}
\begin{minipage}[t]{0.405\textwidth}
\SlideImage{29}{../Materials/2026/3DM_11_Food-Delivery_SS26.pdf}
\begin{adjustbox}{max totalsize={\textwidth}{0.43\textheight},center}
\begin{minipage}{\textwidth}\scriptsize
\NoteSection{日本語訳}\begin{itemize}
\item Online and offline prediction of food handover time（この用語は講義スライド上の専門語として扱う）
\item https://www.infoq.com/fr/articles/uber-eats-time-predictions/（この用語は講義スライド上の専門語として扱う）
\end{itemize}\NoteSection{試験対策ポイント}\begin{itemize}
\item Platform-based food deliveryは、需要・供給・推薦・配送の複数意思決定が連鎖する問題である。
\item food recommendationとdelivery optimizationを分けつつ、相互作用を説明する。
\end{itemize}
\end{minipage}
\end{adjustbox}
\end{minipage}\hfill
\begin{minipage}[t]{0.58\textwidth}
\begin{adjustbox}{max totalsize={\textwidth}{0.89\textheight},center}
\begin{minipage}{\textwidth}\scriptsize
\NoteSection{解説}このページでは「Online and offline prediction of food handover time」を理解する。スライド上の表現は短いが、試験では定義、具体例、モデル上の意味、手法の限界まで説明できる必要がある。\par
Platform-based food delivery は、顧客、レストラン、配達員、プラットフォームが相互作用する多面的なサービスである。顧客は料理を選び、レストランは注文を調理し、配達員は複数注文を運び、プラットフォームは推薦、価格、割当、ルーティングを決める。\par
この問題では、単に最短ルートを作るだけでは不十分である。推薦システムが特定店舗に需要を集中させると、調理待ちや配達遅延が増える可能性がある。逆に、配送能力を考慮した推薦を行えば、需要を空いている店舗や配送しやすい地域へ誘導できる。\par
food recommendation では、content-based recommendation, collaborative filtering, ontology graph, diversificationなどが扱われる。顧客の好みだけでなく、料理の種類、価格、距離、調理時間、配達可能性などが推薦品質に影響する。\par
配送最適化では、注文受諾、配達員割当、ピックアップ順、ドロップオフ順、複数注文のバッチングが問題になる。時間窓、料理の品質劣化、顧客待ち時間、配達員の公平性なども考慮する必要がある。\par
試験では、食品配送をMDP/ADPとして設計する場合、状態に未処理注文、配達員位置、レストラン調理状態、需要予測を入れ、決定に推薦・注文受諾・割当・ルーティングを入れる。報酬には収益、遅延ペナルティ、顧客満足度、配達員稼働を含めるとよい。\par\NoteSection{確認問題と解答}\begin{enumerate}\item \textbf{L11-S029: 「Online and offline prediction of food handover time」の概念を、定義・具体例・モデル上の意味の順に説明せよ。}\\定義としては、Platform-based food deliveryは、需要・供給・推薦・配送の複数意思決定が連鎖する問題である。。具体例では、配送・在庫・フードデリバリーのような現実問題を用い、状態、決定、外生情報、報酬、または情報モデル/意思決定モデルのどこに対応するかを明示する。モデル上の意味として、この概念が目的関数、制約、状態表現、将来価値、または方策選択にどのような影響を与えるかを書く。試験では、用語だけを列挙せず、入力データから意思決定、さらに現実の実装へつながる流れを文章で示すことが重要である。\item \textbf{L11-S029: このスライドの考え方を、講義全体のDDDMプロセスの中に位置づけよ。}\\DDDMプロセスでは、現実世界のデータを集約して情報モデルを作り、意思決定モデルまたは方策で行動を選び、実装後に結果を観測して次の状態へ進む。このスライドの概念は、そのどこか一部だけではなく、データ表現、意思決定、将来不確実性、計算可能性のいずれかに関わる。答案では、まずこの概念がどの段階に属するかを述べ、次に他の段階へどう影響するかを書く。例えば価値関数なら将来価値を現在決定へ戻す役割、FIFOなら旅行時間情報モデルの整合性を保証する役割、CFAなら目的関数を調整して将来に良い行動を誘導する役割である。\end{enumerate}
\end{minipage}
\end{adjustbox}
\end{minipage}


\clearpage
\SlideHeader{Lecture 11 - Platform-based Food Delivery}{030}{Data acquisition, training, evaluation and deployment}
\begin{minipage}[t]{0.405\textwidth}
\SlideImage{30}{../Materials/2026/3DM_11_Food-Delivery_SS26.pdf}
\begin{adjustbox}{max totalsize={\textwidth}{0.43\textheight},center}
\begin{minipage}{\textwidth}\scriptsize
\NoteSection{日本語訳}\begin{itemize}
\item Data acquisition, training, evaluation and deployment（この用語は講義スライド上の専門語として扱う）
\item https://www.infoq.com/fr/articles/uber-eats-time-predictions/（この用語は講義スライド上の専門語として扱う）
\end{itemize}\NoteSection{試験対策ポイント}\begin{itemize}
\item Platform-based food deliveryは、需要・供給・推薦・配送の複数意思決定が連鎖する問題である。
\item food recommendationとdelivery optimizationを分けつつ、相互作用を説明する。
\end{itemize}
\end{minipage}
\end{adjustbox}
\end{minipage}\hfill
\begin{minipage}[t]{0.58\textwidth}
\begin{adjustbox}{max totalsize={\textwidth}{0.89\textheight},center}
\begin{minipage}{\textwidth}\scriptsize
\NoteSection{解説}このページでは「Data acquisition, training, evaluation and deployment」を理解する。スライド上の表現は短いが、試験では定義、具体例、モデル上の意味、手法の限界まで説明できる必要がある。\par
Platform-based food delivery は、顧客、レストラン、配達員、プラットフォームが相互作用する多面的なサービスである。顧客は料理を選び、レストランは注文を調理し、配達員は複数注文を運び、プラットフォームは推薦、価格、割当、ルーティングを決める。\par
この問題では、単に最短ルートを作るだけでは不十分である。推薦システムが特定店舗に需要を集中させると、調理待ちや配達遅延が増える可能性がある。逆に、配送能力を考慮した推薦を行えば、需要を空いている店舗や配送しやすい地域へ誘導できる。\par
food recommendation では、content-based recommendation, collaborative filtering, ontology graph, diversificationなどが扱われる。顧客の好みだけでなく、料理の種類、価格、距離、調理時間、配達可能性などが推薦品質に影響する。\par
配送最適化では、注文受諾、配達員割当、ピックアップ順、ドロップオフ順、複数注文のバッチングが問題になる。時間窓、料理の品質劣化、顧客待ち時間、配達員の公平性なども考慮する必要がある。\par
試験では、食品配送をMDP/ADPとして設計する場合、状態に未処理注文、配達員位置、レストラン調理状態、需要予測を入れ、決定に推薦・注文受諾・割当・ルーティングを入れる。報酬には収益、遅延ペナルティ、顧客満足度、配達員稼働を含めるとよい。\par\NoteSection{確認問題と解答}\begin{enumerate}\item \textbf{L11-S030: 「Data acquisition, training, evaluation and deployment」の概念を、定義・具体例・モデル上の意味の順に説明せよ。}\\定義としては、Platform-based food deliveryは、需要・供給・推薦・配送の複数意思決定が連鎖する問題である。。具体例では、配送・在庫・フードデリバリーのような現実問題を用い、状態、決定、外生情報、報酬、または情報モデル/意思決定モデルのどこに対応するかを明示する。モデル上の意味として、この概念が目的関数、制約、状態表現、将来価値、または方策選択にどのような影響を与えるかを書く。試験では、用語だけを列挙せず、入力データから意思決定、さらに現実の実装へつながる流れを文章で示すことが重要である。\item \textbf{L11-S030: このスライドの考え方を、講義全体のDDDMプロセスの中に位置づけよ。}\\DDDMプロセスでは、現実世界のデータを集約して情報モデルを作り、意思決定モデルまたは方策で行動を選び、実装後に結果を観測して次の状態へ進む。このスライドの概念は、そのどこか一部だけではなく、データ表現、意思決定、将来不確実性、計算可能性のいずれかに関わる。答案では、まずこの概念がどの段階に属するかを述べ、次に他の段階へどう影響するかを書く。例えば価値関数なら将来価値を現在決定へ戻す役割、FIFOなら旅行時間情報モデルの整合性を保証する役割、CFAなら目的関数を調整して将来に良い行動を誘導する役割である。\end{enumerate}
\end{minipage}
\end{adjustbox}
\end{minipage}


\clearpage
\SlideHeader{Lecture 11 - Platform-based Food Delivery}{031}{Analysis of driver‘s activities for selected restaurants}
\begin{minipage}[t]{0.405\textwidth}
\SlideImage{31}{../Materials/2026/3DM_11_Food-Delivery_SS26.pdf}
\begin{adjustbox}{max totalsize={\textwidth}{0.43\textheight},center}
\begin{minipage}{\textwidth}\scriptsize
\NoteSection{日本語訳}\begin{itemize}
\item Analysis of driver‘s activities for selected restaurants（この用語は講義スライド上の専門語として扱う）
\item https://www.infoq.com/fr/articles/uber-eats-time-predictions/（この用語は講義スライド上の専門語として扱う）
\item Restaurants are ranked in accordance to duration and variation of food handover time（この用語は講義スライド上の専門語として扱う）
\end{itemize}\NoteSection{試験対策ポイント}\begin{itemize}
\item Platform-based food deliveryは、需要・供給・推薦・配送の複数意思決定が連鎖する問題である。
\item food recommendationとdelivery optimizationを分けつつ、相互作用を説明する。
\end{itemize}
\end{minipage}
\end{adjustbox}
\end{minipage}\hfill
\begin{minipage}[t]{0.58\textwidth}
\begin{adjustbox}{max totalsize={\textwidth}{0.89\textheight},center}
\begin{minipage}{\textwidth}\scriptsize
\NoteSection{解説}このページでは「Analysis of driver‘s activities for selected restaurants」を理解する。スライド上の表現は短いが、試験では定義、具体例、モデル上の意味、手法の限界まで説明できる必要がある。\par
Platform-based food delivery は、顧客、レストラン、配達員、プラットフォームが相互作用する多面的なサービスである。顧客は料理を選び、レストランは注文を調理し、配達員は複数注文を運び、プラットフォームは推薦、価格、割当、ルーティングを決める。\par
この問題では、単に最短ルートを作るだけでは不十分である。推薦システムが特定店舗に需要を集中させると、調理待ちや配達遅延が増える可能性がある。逆に、配送能力を考慮した推薦を行えば、需要を空いている店舗や配送しやすい地域へ誘導できる。\par
food recommendation では、content-based recommendation, collaborative filtering, ontology graph, diversificationなどが扱われる。顧客の好みだけでなく、料理の種類、価格、距離、調理時間、配達可能性などが推薦品質に影響する。\par
配送最適化では、注文受諾、配達員割当、ピックアップ順、ドロップオフ順、複数注文のバッチングが問題になる。時間窓、料理の品質劣化、顧客待ち時間、配達員の公平性なども考慮する必要がある。\par
試験では、食品配送をMDP/ADPとして設計する場合、状態に未処理注文、配達員位置、レストラン調理状態、需要予測を入れ、決定に推薦・注文受諾・割当・ルーティングを入れる。報酬には収益、遅延ペナルティ、顧客満足度、配達員稼働を含めるとよい。\par\NoteSection{確認問題と解答}\begin{enumerate}\item \textbf{L11-S031: 「Analysis of driver‘s activities for selected restaurants」の概念を、定義・具体例・モデル上の意味の順に説明せよ。}\\定義としては、Platform-based food deliveryは、需要・供給・推薦・配送の複数意思決定が連鎖する問題である。。具体例では、配送・在庫・フードデリバリーのような現実問題を用い、状態、決定、外生情報、報酬、または情報モデル/意思決定モデルのどこに対応するかを明示する。モデル上の意味として、この概念が目的関数、制約、状態表現、将来価値、または方策選択にどのような影響を与えるかを書く。試験では、用語だけを列挙せず、入力データから意思決定、さらに現実の実装へつながる流れを文章で示すことが重要である。\item \textbf{L11-S031: このスライドの考え方を、講義全体のDDDMプロセスの中に位置づけよ。}\\DDDMプロセスでは、現実世界のデータを集約して情報モデルを作り、意思決定モデルまたは方策で行動を選び、実装後に結果を観測して次の状態へ進む。このスライドの概念は、そのどこか一部だけではなく、データ表現、意思決定、将来不確実性、計算可能性のいずれかに関わる。答案では、まずこの概念がどの段階に属するかを述べ、次に他の段階へどう影響するかを書く。例えば価値関数なら将来価値を現在決定へ戻す役割、FIFOなら旅行時間情報モデルの整合性を保証する役割、CFAなら目的関数を調整して将来に良い行動を誘導する役割である。\end{enumerate}
\end{minipage}
\end{adjustbox}
\end{minipage}


\clearpage
\SlideHeader{Lecture 11 - Platform-based Food Delivery}{032}{Prediction of food handover at restaurant}
\begin{minipage}[t]{0.405\textwidth}
\SlideImage{32}{../Materials/2026/3DM_11_Food-Delivery_SS26.pdf}
\begin{adjustbox}{max totalsize={\textwidth}{0.43\textheight},center}
\begin{minipage}{\textwidth}\scriptsize
\NoteSection{日本語訳}\begin{itemize}
\item Prediction of food handover at restaurant（この用語は講義スライド上の専門語として扱う）
\item ▪ Predictive modeling task, 教師あり学習
\item ▪ Inputs like restaurant, time of day, \#orders placed, size of order,（この用語は講義スライド上の専門語として扱う）
\item driver, dispatch time（この用語は講義スライド上の専門語として扱う）
\item ▪ Predicted output is the time span between driver arrival and（この用語は講義スライド上の専門語として扱う）
\item handover of food（この用語は講義スライド上の専門語として扱う）
\item ▪ Phases of the process depend on each other → Markov model（この用語は講義スライド上の専門語として扱う）
\item ▪ Prediction e.g. by recurrent neural networks, that couple ANN（この用語は講義スライド上の専門語として扱う）
\item https://de.wikipedia.org/wiki/Rekurrentes\_neuronales\_Netz\#/me（この用語は講義スライド上の専門語として扱う）
\end{itemize}\NoteSection{試験対策ポイント}\begin{itemize}
\item Platform-based food deliveryは、需要・供給・推薦・配送の複数意思決定が連鎖する問題である。
\item food recommendationとdelivery optimizationを分けつつ、相互作用を説明する。
\end{itemize}
\end{minipage}
\end{adjustbox}
\end{minipage}\hfill
\begin{minipage}[t]{0.58\textwidth}
\begin{adjustbox}{max totalsize={\textwidth}{0.89\textheight},center}
\begin{minipage}{\textwidth}\scriptsize
\NoteSection{解説}このページでは「Prediction of food handover at restaurant」を理解する。スライド上の表現は短いが、試験では定義、具体例、モデル上の意味、手法の限界まで説明できる必要がある。\par
Platform-based food delivery は、顧客、レストラン、配達員、プラットフォームが相互作用する多面的なサービスである。顧客は料理を選び、レストランは注文を調理し、配達員は複数注文を運び、プラットフォームは推薦、価格、割当、ルーティングを決める。\par
この問題では、単に最短ルートを作るだけでは不十分である。推薦システムが特定店舗に需要を集中させると、調理待ちや配達遅延が増える可能性がある。逆に、配送能力を考慮した推薦を行えば、需要を空いている店舗や配送しやすい地域へ誘導できる。\par
food recommendation では、content-based recommendation, collaborative filtering, ontology graph, diversificationなどが扱われる。顧客の好みだけでなく、料理の種類、価格、距離、調理時間、配達可能性などが推薦品質に影響する。\par
配送最適化では、注文受諾、配達員割当、ピックアップ順、ドロップオフ順、複数注文のバッチングが問題になる。時間窓、料理の品質劣化、顧客待ち時間、配達員の公平性なども考慮する必要がある。\par
試験では、食品配送をMDP/ADPとして設計する場合、状態に未処理注文、配達員位置、レストラン調理状態、需要予測を入れ、決定に推薦・注文受諾・割当・ルーティングを入れる。報酬には収益、遅延ペナルティ、顧客満足度、配達員稼働を含めるとよい。\par\NoteSection{確認問題と解答}\begin{enumerate}\item \textbf{L11-S032: 「Prediction of food handover at restaurant」の概念を、定義・具体例・モデル上の意味の順に説明せよ。}\\定義としては、Platform-based food deliveryは、需要・供給・推薦・配送の複数意思決定が連鎖する問題である。。具体例では、配送・在庫・フードデリバリーのような現実問題を用い、状態、決定、外生情報、報酬、または情報モデル/意思決定モデルのどこに対応するかを明示する。モデル上の意味として、この概念が目的関数、制約、状態表現、将来価値、または方策選択にどのような影響を与えるかを書く。試験では、用語だけを列挙せず、入力データから意思決定、さらに現実の実装へつながる流れを文章で示すことが重要である。\item \textbf{L11-S032: このスライドの考え方を、講義全体のDDDMプロセスの中に位置づけよ。}\\DDDMプロセスでは、現実世界のデータを集約して情報モデルを作り、意思決定モデルまたは方策で行動を選び、実装後に結果を観測して次の状態へ進む。このスライドの概念は、そのどこか一部だけではなく、データ表現、意思決定、将来不確実性、計算可能性のいずれかに関わる。答案では、まずこの概念がどの段階に属するかを述べ、次に他の段階へどう影響するかを書く。例えば価値関数なら将来価値を現在決定へ戻す役割、FIFOなら旅行時間情報モデルの整合性を保証する役割、CFAなら目的関数を調整して将来に良い行動を誘導する役割である。\end{enumerate}
\end{minipage}
\end{adjustbox}
\end{minipage}


\clearpage
\SlideHeader{Lecture 11 - Platform-based Food Delivery}{033}{Adjusting prediction by real world observation}
\begin{minipage}[t]{0.405\textwidth}
\SlideImage{33}{../Materials/2026/3DM_11_Food-Delivery_SS26.pdf}
\begin{adjustbox}{max totalsize={\textwidth}{0.43\textheight},center}
\begin{minipage}{\textwidth}\scriptsize
\NoteSection{日本語訳}\begin{itemize}
\item Adjusting prediction by real world observation（この用語は講義スライド上の専門語として扱う）
\item Activity modes that are recognized by Android OS（この用語は講義スライド上の専門語として扱う）
\item Sequence of activities to be carried out by driver（この用語は講義スライド上の専門語として扱う）
\item https://www.uber.com/en-DE/blog/uber-eats-trip-（この用語は講義スライド上の専門語として扱う）
\item optimization/（この用語は講義スライド上の専門語として扱う）
\end{itemize}\NoteSection{試験対策ポイント}\begin{itemize}
\item Platform-based food deliveryは、需要・供給・推薦・配送の複数意思決定が連鎖する問題である。
\item food recommendationとdelivery optimizationを分けつつ、相互作用を説明する。
\end{itemize}
\end{minipage}
\end{adjustbox}
\end{minipage}\hfill
\begin{minipage}[t]{0.58\textwidth}
\begin{adjustbox}{max totalsize={\textwidth}{0.89\textheight},center}
\begin{minipage}{\textwidth}\scriptsize
\NoteSection{解説}このページでは「Adjusting prediction by real world observation」を理解する。スライド上の表現は短いが、試験では定義、具体例、モデル上の意味、手法の限界まで説明できる必要がある。\par
Platform-based food delivery は、顧客、レストラン、配達員、プラットフォームが相互作用する多面的なサービスである。顧客は料理を選び、レストランは注文を調理し、配達員は複数注文を運び、プラットフォームは推薦、価格、割当、ルーティングを決める。\par
この問題では、単に最短ルートを作るだけでは不十分である。推薦システムが特定店舗に需要を集中させると、調理待ちや配達遅延が増える可能性がある。逆に、配送能力を考慮した推薦を行えば、需要を空いている店舗や配送しやすい地域へ誘導できる。\par
food recommendation では、content-based recommendation, collaborative filtering, ontology graph, diversificationなどが扱われる。顧客の好みだけでなく、料理の種類、価格、距離、調理時間、配達可能性などが推薦品質に影響する。\par
配送最適化では、注文受諾、配達員割当、ピックアップ順、ドロップオフ順、複数注文のバッチングが問題になる。時間窓、料理の品質劣化、顧客待ち時間、配達員の公平性なども考慮する必要がある。\par
試験では、食品配送をMDP/ADPとして設計する場合、状態に未処理注文、配達員位置、レストラン調理状態、需要予測を入れ、決定に推薦・注文受諾・割当・ルーティングを入れる。報酬には収益、遅延ペナルティ、顧客満足度、配達員稼働を含めるとよい。\par\NoteSection{確認問題と解答}\begin{enumerate}\item \textbf{L11-S033: 「Adjusting prediction by real world observation」の概念を、定義・具体例・モデル上の意味の順に説明せよ。}\\定義としては、Platform-based food deliveryは、需要・供給・推薦・配送の複数意思決定が連鎖する問題である。。具体例では、配送・在庫・フードデリバリーのような現実問題を用い、状態、決定、外生情報、報酬、または情報モデル/意思決定モデルのどこに対応するかを明示する。モデル上の意味として、この概念が目的関数、制約、状態表現、将来価値、または方策選択にどのような影響を与えるかを書く。試験では、用語だけを列挙せず、入力データから意思決定、さらに現実の実装へつながる流れを文章で示すことが重要である。\item \textbf{L11-S033: このスライドの考え方を、講義全体のDDDMプロセスの中に位置づけよ。}\\DDDMプロセスでは、現実世界のデータを集約して情報モデルを作り、意思決定モデルまたは方策で行動を選び、実装後に結果を観測して次の状態へ進む。このスライドの概念は、そのどこか一部だけではなく、データ表現、意思決定、将来不確実性、計算可能性のいずれかに関わる。答案では、まずこの概念がどの段階に属するかを述べ、次に他の段階へどう影響するかを書く。例えば価値関数なら将来価値を現在決定へ戻す役割、FIFOなら旅行時間情報モデルの整合性を保証する役割、CFAなら目的関数を調整して将来に良い行動を誘導する役割である。\end{enumerate}
\end{minipage}
\end{adjustbox}
\end{minipage}


\clearpage
\SlideHeader{Lecture 11 - Platform-based Food Delivery}{034}{From order creation to food drop-off}
\begin{minipage}[t]{0.405\textwidth}
\SlideImage{34}{../Materials/2026/3DM_11_Food-Delivery_SS26.pdf}
\begin{adjustbox}{max totalsize={\textwidth}{0.43\textheight},center}
\begin{minipage}{\textwidth}\scriptsize
\NoteSection{日本語訳}\begin{itemize}
\item From order creation to food drop-off（この用語は講義スライド上の専門語として扱う）
\item https://www.infoq.com/fr/articles/uber-eats-time-predictions/（この用語は講義スライド上の専門語として扱う）
\end{itemize}\NoteSection{試験対策ポイント}\begin{itemize}
\item Platform-based food deliveryは、需要・供給・推薦・配送の複数意思決定が連鎖する問題である。
\item food recommendationとdelivery optimizationを分けつつ、相互作用を説明する。
\end{itemize}
\end{minipage}
\end{adjustbox}
\end{minipage}\hfill
\begin{minipage}[t]{0.58\textwidth}
\begin{adjustbox}{max totalsize={\textwidth}{0.89\textheight},center}
\begin{minipage}{\textwidth}\scriptsize
\NoteSection{解説}このページでは「From order creation to food drop-off」を理解する。スライド上の表現は短いが、試験では定義、具体例、モデル上の意味、手法の限界まで説明できる必要がある。\par
Platform-based food delivery は、顧客、レストラン、配達員、プラットフォームが相互作用する多面的なサービスである。顧客は料理を選び、レストランは注文を調理し、配達員は複数注文を運び、プラットフォームは推薦、価格、割当、ルーティングを決める。\par
この問題では、単に最短ルートを作るだけでは不十分である。推薦システムが特定店舗に需要を集中させると、調理待ちや配達遅延が増える可能性がある。逆に、配送能力を考慮した推薦を行えば、需要を空いている店舗や配送しやすい地域へ誘導できる。\par
food recommendation では、content-based recommendation, collaborative filtering, ontology graph, diversificationなどが扱われる。顧客の好みだけでなく、料理の種類、価格、距離、調理時間、配達可能性などが推薦品質に影響する。\par
配送最適化では、注文受諾、配達員割当、ピックアップ順、ドロップオフ順、複数注文のバッチングが問題になる。時間窓、料理の品質劣化、顧客待ち時間、配達員の公平性なども考慮する必要がある。\par
試験では、食品配送をMDP/ADPとして設計する場合、状態に未処理注文、配達員位置、レストラン調理状態、需要予測を入れ、決定に推薦・注文受諾・割当・ルーティングを入れる。報酬には収益、遅延ペナルティ、顧客満足度、配達員稼働を含めるとよい。\par\NoteSection{確認問題と解答}\begin{enumerate}\item \textbf{L11-S034: 「From order creation to food drop-off」の概念を、定義・具体例・モデル上の意味の順に説明せよ。}\\定義としては、Platform-based food deliveryは、需要・供給・推薦・配送の複数意思決定が連鎖する問題である。。具体例では、配送・在庫・フードデリバリーのような現実問題を用い、状態、決定、外生情報、報酬、または情報モデル/意思決定モデルのどこに対応するかを明示する。モデル上の意味として、この概念が目的関数、制約、状態表現、将来価値、または方策選択にどのような影響を与えるかを書く。試験では、用語だけを列挙せず、入力データから意思決定、さらに現実の実装へつながる流れを文章で示すことが重要である。\item \textbf{L11-S034: このスライドの考え方を、講義全体のDDDMプロセスの中に位置づけよ。}\\DDDMプロセスでは、現実世界のデータを集約して情報モデルを作り、意思決定モデルまたは方策で行動を選び、実装後に結果を観測して次の状態へ進む。このスライドの概念は、そのどこか一部だけではなく、データ表現、意思決定、将来不確実性、計算可能性のいずれかに関わる。答案では、まずこの概念がどの段階に属するかを述べ、次に他の段階へどう影響するかを書く。例えば価値関数なら将来価値を現在決定へ戻す役割、FIFOなら旅行時間情報モデルの整合性を保証する役割、CFAなら目的関数を調整して将来に良い行動を誘導する役割である。\end{enumerate}
\end{minipage}
\end{adjustbox}
\end{minipage}


\clearpage
\SlideHeader{Lecture 11 - Platform-based Food Delivery}{035}{Summary}
\begin{minipage}[t]{0.405\textwidth}
\SlideImage{35}{../Materials/2026/3DM_11_Food-Delivery_SS26.pdf}
\begin{adjustbox}{max totalsize={\textwidth}{0.43\textheight},center}
\begin{minipage}{\textwidth}\scriptsize
\NoteSection{日本語訳}\begin{itemize}
\item まとめ
\item ▪ Platform business models require online automated 決定 making
\item ▪ 決定 making incorporates prediction by means of data analytics
\item ▪ Online optimization ensures seamless transaction of meal delivery（この用語は講義スライド上の専門語として扱う）
\item ▪ Expected realization of individual objectives determine food 推薦
\item ▪ Realized 推薦s determine conversion rate → to be 最大化するd
\end{itemize}\NoteSection{試験対策ポイント}\begin{itemize}
\item Platform-based food deliveryは、需要・供給・推薦・配送の複数意思決定が連鎖する問題である。
\item food recommendationとdelivery optimizationを分けつつ、相互作用を説明する。
\end{itemize}
\end{minipage}
\end{adjustbox}
\end{minipage}\hfill
\begin{minipage}[t]{0.58\textwidth}
\begin{adjustbox}{max totalsize={\textwidth}{0.89\textheight},center}
\begin{minipage}{\textwidth}\scriptsize
\NoteSection{解説}このページでは「Summary」を理解する。スライド上の表現は短いが、試験では定義、具体例、モデル上の意味、手法の限界まで説明できる必要がある。\par
Platform-based food delivery は、顧客、レストラン、配達員、プラットフォームが相互作用する多面的なサービスである。顧客は料理を選び、レストランは注文を調理し、配達員は複数注文を運び、プラットフォームは推薦、価格、割当、ルーティングを決める。\par
この問題では、単に最短ルートを作るだけでは不十分である。推薦システムが特定店舗に需要を集中させると、調理待ちや配達遅延が増える可能性がある。逆に、配送能力を考慮した推薦を行えば、需要を空いている店舗や配送しやすい地域へ誘導できる。\par
food recommendation では、content-based recommendation, collaborative filtering, ontology graph, diversificationなどが扱われる。顧客の好みだけでなく、料理の種類、価格、距離、調理時間、配達可能性などが推薦品質に影響する。\par
配送最適化では、注文受諾、配達員割当、ピックアップ順、ドロップオフ順、複数注文のバッチングが問題になる。時間窓、料理の品質劣化、顧客待ち時間、配達員の公平性なども考慮する必要がある。\par
試験では、食品配送をMDP/ADPとして設計する場合、状態に未処理注文、配達員位置、レストラン調理状態、需要予測を入れ、決定に推薦・注文受諾・割当・ルーティングを入れる。報酬には収益、遅延ペナルティ、顧客満足度、配達員稼働を含めるとよい。\par\NoteSection{確認問題と解答}\begin{enumerate}\item \textbf{L11-S035: 「Summary」の概念を、定義・具体例・モデル上の意味の順に説明せよ。}\\定義としては、Platform-based food deliveryは、需要・供給・推薦・配送の複数意思決定が連鎖する問題である。。具体例では、配送・在庫・フードデリバリーのような現実問題を用い、状態、決定、外生情報、報酬、または情報モデル/意思決定モデルのどこに対応するかを明示する。モデル上の意味として、この概念が目的関数、制約、状態表現、将来価値、または方策選択にどのような影響を与えるかを書く。試験では、用語だけを列挙せず、入力データから意思決定、さらに現実の実装へつながる流れを文章で示すことが重要である。\item \textbf{L11-S035: このスライドの考え方を、講義全体のDDDMプロセスの中に位置づけよ。}\\DDDMプロセスでは、現実世界のデータを集約して情報モデルを作り、意思決定モデルまたは方策で行動を選び、実装後に結果を観測して次の状態へ進む。このスライドの概念は、そのどこか一部だけではなく、データ表現、意思決定、将来不確実性、計算可能性のいずれかに関わる。答案では、まずこの概念がどの段階に属するかを述べ、次に他の段階へどう影響するかを書く。例えば価値関数なら将来価値を現在決定へ戻す役割、FIFOなら旅行時間情報モデルの整合性を保証する役割、CFAなら目的関数を調整して将来に良い行動を誘導する役割である。\end{enumerate}
\end{minipage}
\end{adjustbox}
\end{minipage}


\clearpage
\SlideHeader{Lecture 11 - Platform-based Food Delivery}{036}{Literature}
\begin{minipage}[t]{0.405\textwidth}
\SlideImage{36}{../Materials/2026/3DM_11_Food-Delivery_SS26.pdf}
\begin{adjustbox}{max totalsize={\textwidth}{0.43\textheight},center}
\begin{minipage}{\textwidth}\scriptsize
\NoteSection{日本語訳}\begin{itemize}
\item 文献
\item ▪ Arriving Now: The New Uber Eats (2020)（この用語は講義スライド上の専門語として扱う）
\item https://www.uber.com/newsroom/arriving-now-the-new-uber-eats/（この用語は講義スライド上の専門語として扱う）
\item ▪ Food Discovery with Uber Eats: Recommending for the Marketplace (2018)（この用語は講義スライド上の専門語として扱う）
\item https://www.uber.com/en-DE/blog/uber-eats-recommending-marketplace/（この用語は講義スライド上の専門語として扱う）
\item ▪ Food Discovery with Uber Eats: Building a Query Understanding Engine (2018)（この用語は講義スライド上の専門語として扱う）
\item https://www.uber.com/en-BE/blog/uber-eats-query-understanding/（この用語は講義スライド上の専門語として扱う）
\item ▪ How Does Uber Eats Assign Drivers? (2024)（この用語は講義スライド上の専門語として扱う）
\item https://www.withpara.com/blog/how-does-uber-eats-assign-drivers（この用語は講義スライド上の専門語として扱う）
\end{itemize}
\end{minipage}
\end{adjustbox}
\end{minipage}\hfill
\begin{minipage}[t]{0.58\textwidth}
\begin{adjustbox}{max totalsize={\textwidth}{0.89\textheight},center}
\begin{minipage}{\textwidth}\scriptsize
\NoteSection{注記}このページは表紙・目次・事務情報・導入画像が中心であるため、無理に確認問題や過去問を付けない。
\end{minipage}
\end{adjustbox}
\end{minipage}

\clearpage\ExamHeader{Lecture 11 - Platform-based Food Delivery}
\ExamQuestion{SS23 Aufgabe 5 (8 点)}
\ExamSubsection{問題内容}
レストラン配送のCFA/PFA。設問では、遅延リスクを費用近似または方策近似で扱い、注文統合の判断を説明すること。 ことが求められる。\par
\ExamSubsection{満点答案}
満点答案では、Real Worldから観測データを集約してInformation Modelを作り、そこからDecision Modelを作り、解をImplementation/Disaggregationで現実の行動へ戻す流れを説明する。Information Modelは現実の圧縮表現であり、Decision Modelは決定変数、目的関数、制約により行動選択を評価する構造である。\par
\ExamSubsection{具体例による解説}
具体例として、配送問題なら、顧客注文・車両位置・旅行時間を情報モデルにし、訪問順や割当を意思決定モデルで決める。在庫問題なら、在庫水準を状態、注文量を決定、需要を外生情報、在庫更新式を遷移、欠品費・保管費を費用として整理する。問題文の文脈に合わせて、抽象用語を必ず具体的な変数や行動に置き換える。\par
\ExamSubsection{採点で落とさない要素}
\begin{itemize}
\item 設問が求める概念の定義を最初に明確に書く。
\item 講義の専門語を列挙するだけでなく、現実例とモデル要素へ対応付ける。
\item 利点だけでなく限界・注意点も最低一つ述べる。
\item 式が出る問題では、記号の意味を文章で説明する。
\item 新しい事例問題では、状態・決定・外生情報・遷移・報酬/費用の順に整理する。
\end{itemize}
\vspace{2mm}\hrule\vspace{2mm}